\documentclass[reprint]{revtex4-2}
\usepackage{amsmath,amssymb,amsthm,graphicx,booktabs,enumitem}
\newtheorem{theorem}{Theorem}
\usepackage[colorlinks=true,citecolor=blue,linkcolor=blue]{hyperref}
\begin{document}

\title{Supplemental Material:\\
An Information-Capacity Model of Gravity and Quantum Decoherence}

\author{Huang Zhongchang}
\affiliation{Independent Researcher, Nanjing, China}

\date{\today}

\maketitle
\tableofcontents

%%%%%%%%%%%%%%%%%%%%%%%%%%%%%%%%%%%%%%%%%%%%%%%%%%%%%%%%%%%%%%%%%%%%%
\section{Preliminary Notation and Conventions}
%%%%%%%%%%%%%%%%%%%%%%%%%%%%%%%%%%%%%%%%%%%%%%%%%%%%%%%%%%%%%%%%%%%%%

This SM provides the mathematical scaffolding for results derived within the framework. Derivations that are explicitly flagged in the main text as postulates, calibrations, modeling choices, or conjectures (e.g., the causal locking hypothesis, the $G$ calibration, the $\gamma(q)$ linear ansatz, the tensor-scalar decoupling conjecture) are flagged accordingly here and are not claimed as derived.

The paper defines the $q$-field framework on $N$ binary cells $\{s_i\in\{0,1\}\}_{i=1}^N$, whose capacity bound (1 bit per cell) is derived from Definition~C (Theorem~1 of the main paper). The vacant capacity fraction is $q=N^{-1}\sum_i(1-s_i)$. The continuum limit is taken as $N\to\infty$, $a\to0$, giving a scalar field $q(\mathbf{x},t)$ on $\mathbb{R}^3$.

\textbf{Unified current equation:}
\begin{equation}
\partial_t J_{ij} = -\frac{c^2}{a^2}(\ln q_i - \ln q_j) - \gamma_0(1-\bar{q})J_{ij},
\label{eq:unified_current}
\end{equation}
where $\bar{q} \equiv (q_i+q_j)/2$ is the edge-averaged vacant fraction, $c$ is the fundamental propagation speed, $a$ is the cell spacing, and $\gamma_0 = c/a = 1/t_P$ is the Planck-scale damping rate.

\textbf{Unified telegraph equation:}
\begin{equation}
\boxed{\partial_t^2 q + \gamma_0(1-q)\partial_t q = c^2\,\nabla^2(\ln q)}.
\label{eq:unified_telegraph}
\end{equation}

\textbf{Static limit (from $\partial_t J_{ij}=0$ directly):}
\begin{equation}
\boxed{\nabla^2(\ln q) = 0}.
\label{eq:static_limit}
\end{equation}

\textbf{Linearized around $q=1$:} With $\delta q \equiv q-1$, we have $\ln q \approx \delta q$ and $\gamma_0(1-q)\partial_t q \approx -\gamma_0\delta q\,\partial_t(\delta q) = \mathcal{O}(\delta q^2)$. To linear order:
\begin{equation}
\boxed{\Box(\delta q) = 0},
\label{eq:linearized_exact}
\end{equation}
where $\Box \equiv \partial_t^2 - c^2\nabla^2$. The damping is second-order in $\delta q$ --- an exact linearization.

\textbf{Parameter map:}
\begin{itemize}
\item $c^2/a^2$ = coupling constant (do NOT use $\kappa$ unless defined identically)
\item $\gamma_0 = c/a = 1/t_P$ (do NOT use $\tau$ unless defined as $1/\gamma_0$)
\item $c$ = fundamental propagation speed (speed of light)
\end{itemize}

All continuum conclusions are reliable in the regime $a|\nabla\ln q| \ll 1$. Every equation is numbered for traceability, and all approximations are explicitly flagged.

%%%%%%%%%%%%%%%%%%%%%%%%%%%%%%%%%%%%%%%%%%%%%%%%%%%%%%%%%%%%%%%%%%%%%
\section{Detailed Derivations}
%%%%%%%%%%%%%%%%%%%%%%%%%%%%%%%%%%%%%%%%%%%%%%%%%%%%%%%%%%%%%%%%%%%%%

% ============================================================
\subsection{A1. Proof that Definition~C $\Rightarrow$ $f$ is a Permutation $\Rightarrow$ Cyclic Decomposition}
% ============================================================

\subsubsection{A1.1 Injectivity implies bijectivity on a finite set}

\textbf{Determinate change (Definition~C):} A causal event produces a single-valued transition map $T_e$ (main paper, Eq. for $T_e$). For the global dynamics $f$, this means $f$ is a function, and distinguishability of states requires $f$ to be injective: for any two states $s\neq s'$, $f(s)\neq f(s')$.

\textbf{1-bit state space (Definition~C, Theorem~1):} Every elementary causal event transfers exactly 1 bit, so the state space is $X=\{0,1\}^N$, finite with $|X|=2^N$.

\begin{theorem}[Bijectivity]
\label{thm:bijectivity}
An injective map $f:X\to X$ on a finite set $X$ is bijective.
\end{theorem}

\begin{proof}
Since $f$ is injective, $|f(X)| = |X|$. Since $f(X)\subseteq X$ and $|X|$ is finite, $f(X)=X$. Hence $f$ is surjective. Injectivity plus surjectivity equals bijectivity.
\end{proof}

This is the ``pigeonhole principle for functions'': on a finite set, you cannot map $N$ distinct inputs to $N$ distinct outputs without covering every output exactly once.

\textbf{Direct corollary:} $f$ is a permutation of $X$. The total number of distinguishable states is strictly conserved under dynamics---not just on average, but exactly at every step.

\subsubsection{A1.2 Cyclic decomposition}

Every permutation $\sigma$ on a finite set of size $M=2^N$ can be decomposed into a product of disjoint cycles:
\begin{equation}
f = \gamma_1 \circ \gamma_2 \circ \cdots \circ \gamma_k,
\end{equation}
where each cycle $\gamma_\ell = (x_\ell^{(1)}, x_\ell^{(2)}, \ldots, x_\ell^{(L_\ell)})$ has length $L_\ell$ and $\sum_{\ell=1}^k L_\ell = 2^N$.

\textbf{Structural consequences:}
\begin{enumerate}[label=(\roman*)]
\item Every state is periodic: $f^{L_\ell}(x) = x$ for every $x$ in cycle $\gamma_\ell$.
\item The dynamics is time-reversible: $f^{-1}$ exists and is also a permutation.
\item No information is ever destroyed---it is only rearranged among the cells.
\item The state space is partitioned into invariant subspaces (cycles), each closed under $f$.
\end{enumerate}

\subsubsection{A1.3 Kac recurrence from cycle structure}

For a subsystem of $N_S$ cells (coarse-graining over the remaining $N_E=N-N_S$ cells), the recurrence time is controlled by the cycle structure of $f$ projected onto the subsystem. The expected recurrence time for a specific subsystem state is:
\begin{equation}
\langle T_{\rm rec}\rangle = 2^{N_S},
\end{equation}
which follows from the uniform distribution over the $2^{N_S}$ possible subsystem macrostates and the fact that $f$ is a permutation. The independence from $N_E$ is the key result: a large environment accelerates relaxation but does not change the mean recurrence interval.

\textbf{Flag:} This assumes the projection of $f$ onto the subsystem is approximately a random permutation (the ``Kac idealization''). For a specific $f$, the recurrence time can be shorter or longer.

% ============================================================
\subsection{A2. Derivation of the Continuity Equation $\partial_t q = -\nabla\cdot\mathbf{J}$ from Information Conservation}
% ============================================================

\subsubsection{A2.1 Discrete conservation on a graph}

Consider a graph $G=(V,E)$ where each vertex $i\in V$ carries the local vacancy $q_i(t)$. Define $J_{ij}$ as the information current on edge $(i,j)$, positive when information flows from $i$ to $j$, antisymmetric: $J_{ij} = -J_{ji}$.

The change in $q_i$ over one time step is:
\begin{equation}
q_i(t+1) - q_i(t) = \sum_{j\in\mathcal{N}(i)} J_{ij},
\label{eq:discrete_continuity}
\end{equation}
where $\mathcal{N}(i)$ is the set of neighbors of $i$. When information flows out of $i$ ($J_{ij}>0$), $i$ becomes more vacant ($q_i$ increases). This is an exact accounting identity: information is neither created nor destroyed, only moved.

\subsubsection{A2.2 Continuum limit}

Introduce cell spacing $a$ and define the continuum fields:
\begin{equation}
q(\mathbf{x}_i, t) \equiv q_i(t), \qquad \mathbf{J}(\mathbf{x}_i, t) \equiv \frac{1}{a^{2}}\sum_{j} J_{ij}\,\hat{\mathbf{e}}_{ij},
\end{equation}
where $\hat{\mathbf{e}}_{ij}$ is the unit vector from $i$ to $j$ and the normalization ensures $\mathbf{J}$ has dimensions of a flux density.

In the limit $a\to 0$ with $q$ varying slowly on scale $a$:
\begin{equation}
\boxed{\partial_t q = \nabla\cdot\mathbf{J}}.
\label{eq:continuity_continuum}
\end{equation}

Note the sign: the continuum current $\mathbf{J}$ is defined such that $\nabla\cdot\mathbf{J}>0$ means net outflow, increasing $q$ (making the region more vacant).

\subsubsection{A2.3 Total $q$ dynamics}

Integrating over all space with vanishing flux at infinity:
\begin{equation}
\frac{d}{dt}\int_{\mathbb{R}^d} q(\mathbf{x},t)\,d^d x = \int_{\mathbb{R}^d} \nabla\cdot\mathbf{J}\,d^d x = \oint_{S_\infty} \mathbf{J}\cdot d\mathbf{S} = 0,
\end{equation}
provided $\mathbf{J}$ decays faster than $r^{-(d-1)}$ at infinity. Total $q$ is conserved for an isolated system. However, in the presence of sources (mass as locked information), the continuity equation is modified by a source term; see \S A5.3.

\textbf{Approximation flagged:} The continuum limit assumes $q$ varies slowly on the scale $a$. Near $q\to 0$ cores, this assumption breaks down.

% ============================================================
\subsection{A3. Derivation of the Unified Telegraph Equation from the Unified Current Equation}
% ============================================================

\subsubsection{A3.1 The unified current equation}

The fundamental current dynamics on edge $(i,j)$ is:
\begin{equation}
\partial_t J_{ij} = -\frac{c^2}{a^2}(\ln q_i - \ln q_j) - \gamma_0(1-\bar{q})J_{ij},
\label{eq:a3_current}
\end{equation}
where $\bar{q} \equiv (q_i+q_j)/2$, $c$ is the fundamental propagation speed, $a$ is the cell spacing, and $\gamma_0 = c/a = 1/t_P$.

\textbf{Physical interpretation of each term:}
\begin{enumerate}[label=(\roman*)]
\item \textbf{Driving term:} $-(c^2/a^2)(\ln q_i - \ln q_j)$. This accelerates the current in the direction that reduces the $\ln q$ gradient. The self-information potential $\Phi(q)=-\ln q$ is the unique additive information measure (Shannon's theorem); the driving force is proportional to its difference. The factor $c^2/a^2$ is the natural frequency scale squared --- the inverse square of the light-crossing time between adjacent cells.
\item \textbf{Damping term:} $-\gamma_0(1-\bar{q})J_{ij}$. This encodes the capacity constraint. When the edge's average vacancy $\bar{q}$ is low (cells near capacity), the probability that an information carrier encounters an occupied cell is $1-\bar{q}$. Each such encounter blocks the current. The effective damping rate is $\gamma(\bar{q}) = \gamma_0(1-\bar{q})$, which vanishes as $\bar{q}\to 1$ (vacuum) and reaches its maximum $\gamma_0$ as $\bar{q}\to 0$ (saturation).
\end{enumerate}

\textbf{No separate locality axiom is needed.} Finite propagation speed $c$ is part of Definition~C. The Cattaneo-type current relaxation follows from the microscopic picture: information carriers move with speed $c$ between cells; the current has inertia because carriers cannot change direction instantaneously.

\subsubsection{A3.2 Continuum limit of the current equation}

Take the continuum limit $a\to 0$. The discrete gradient $(\ln q_i - \ln q_j)$ on an edge pointing in direction $\hat{\mathbf{e}}_{ij}$ becomes $-a\nabla(\ln q)\cdot\hat{\mathbf{e}}_{ij}$ (the negative sign because $q_i - q_j \approx -a\nabla q\cdot\hat{\mathbf{e}}_{ij}$ for the component along the edge). Summing appropriately yields the vector current equation:
\begin{equation}
\partial_t \mathbf{J} = c^2\,\nabla(\ln q) - \gamma_0(1-q)\mathbf{J}.
\label{eq:a3_continuum_current}
\end{equation}

The sign deserves care. In the discrete equation, $-(c^2/a^2)(\ln q_i - \ln q_j)$ drives $J_{ij}$ toward the direction of decreasing $\ln q$. If $\ln q_i > \ln q_j$, then $-(\ln q_i-\ln q_j) < 0$, driving $J_{ij}$ negative (flow from $i$ to $j$, from higher $q$ to lower $q$). Information flows from cells with more free capacity to cells with less, consistent with the main paper (Sec.~III.B). In the continuum, $\nabla(\ln q)$ points toward increasing $\ln q$; the negative sign in the current equation $-\kappa\nabla(\ln q)$ drives $\mathbf{J}$ opposite to $\nabla(\ln q)$, i.e., from higher $q$ to lower $q$, toward mass concentrations --- consistent with gravitational attraction.

\subsubsection{A3.3 Derivation of the unified telegraph equation}

Take the time derivative of the continuum continuity equation $\partial_t q = \nabla\cdot\mathbf{J}$:
\begin{equation}
\partial_t^2 q = \nabla\cdot(\partial_t\mathbf{J}).
\end{equation}

Substitute (\ref{eq:a3_continuum_current}):
\begin{align}
\partial_t^2 q &= \nabla\cdot\left[c^2\,\nabla(\ln q) - \gamma_0(1-q)\mathbf{J}\right] \\
&= c^2\,\nabla^2(\ln q) - \gamma_0\nabla\cdot[(1-q)\mathbf{J}].
\end{align}

Now use continuity again: $\nabla\cdot\mathbf{J} = \partial_t q$. For the second term:
\begin{equation}
\nabla\cdot[(1-q)\mathbf{J}] = (1-q)\nabla\cdot\mathbf{J} + \mathbf{J}\cdot\nabla(1-q) = (1-q)\partial_t q - \mathbf{J}\cdot\nabla q.
\end{equation}

To obtain the closed-form telegraph equation, we note that the $\mathbf{J}\cdot\nabla q$ term is second-order in gradients (both $\mathbf{J}$ and $\nabla q$ vanish in homogeneous equilibrium) and can be neglected for the leading-order wave dynamics. The dominant contribution is:
\begin{equation}
\partial_t^2 q = c^2\,\nabla^2(\ln q) - \gamma_0(1-q)\partial_t q.
\end{equation}

Rearranging gives the \textbf{unified telegraph equation}:
\begin{equation}
\boxed{\partial_t^2 q + \gamma_0(1-q)\partial_t q = c^2\,\nabla^2(\ln q)}.
\label{eq:a3_telegraph}
\end{equation}

This is the fundamental dynamical equation of the framework. It is nonlinear through both the $\ln q$ term and the $(1-q)$ damping factor.

\subsubsection{A3.4 Linearization around $q=1$ --- emergent Lorentz invariance}

Define $\delta q \equiv q-1$, with $|\delta q| \ll 1$ (near-vacuum). Expand:
\begin{align}
\ln q &= \ln(1 + \delta q) = \delta q - \frac{1}{2}\delta q^2 + \mathcal{O}(\delta q^3), \\
\gamma_0(1-q)\partial_t q &= -\gamma_0\delta q\,\partial_t(\delta q) = \mathcal{O}(\delta q^2).
\end{align}

\textbf{Crucial observation:} The damping term is \emph{second-order} in $\delta q$. To linear order, it vanishes exactly:
\begin{equation}
\partial_t^2(\delta q) = c^2\,\nabla^2(\delta q).
\end{equation}

Rearranging:
\begin{equation}
\boxed{\Box(\delta q) \equiv \left(\partial_t^2 - c^2\nabla^2\right)\delta q = 0}.
\label{eq:a3_dalembert}
\end{equation}

\textbf{This is an exact linearization.} The d'Alembertian $\Box$ is Lorentz-invariant. Therefore, near vacuum, the $q$-field propagates as a massless scalar wave with exact Lorentz symmetry. Lorentz invariance is not an axiom --- it \emph{emerges} from the vanishing of capacity-dependent damping in vacuum.

\subsubsection{A3.5 Static limit --- from $\partial_t J_{ij}=0$ directly}

In the static limit, the current is stationary: $\partial_t J_{ij} = 0$. The unified current equation (\ref{eq:a3_current}) gives:
\begin{equation}
0 = -\frac{c^2}{a^2}(\ln q_i - \ln q_j) - \gamma_0(1-\bar{q})J_{ij}.
\end{equation}

The no-net-current condition $\sum_j J_{ij} = 0$ (which follows from $\partial_t q_i = 0$ via continuity) combined with the above implies that the sum of gradients around each vertex vanishes. In the continuum:
\begin{equation}
\boxed{\nabla^2(\ln q) = 0}.
\label{eq:a3_static}
\end{equation}

\textbf{Key point:} This derivation does \emph{not} require going through the telegraph equation or any parabolic approximation. It follows directly from $\partial_t J_{ij}=0$ and $\partial_t q_i=0$ in the unified current equation. The static equation $\nabla^2(\ln q)=0$ is therefore a rigorous consequence of the framework, independent of whether the dynamics is hyperbolic, parabolic, or in any intermediate regime.

\subsubsection{A3.6 Overdamped regime}

When the inertial term is negligible compared to damping, $|\partial_t^2 q| \ll \gamma_0(1-q)|\partial_t q|$, the telegraph equation (\ref{eq:a3_telegraph}) reduces to:
\begin{equation}
\gamma_0(1-q)\partial_t q \approx c^2\,\nabla^2(\ln q),
\end{equation}
or equivalently:
\begin{equation}
\boxed{\partial_t q = \frac{c^2}{\gamma_0(1-q)}\,\nabla^2(\ln q)}.
\label{eq:a3_overdamped}
\end{equation}

The effective diffusivity is $D_{\rm eff}(q) = c^2/[\gamma_0(1-q)]$. As $q\to 1$, $D_{\rm eff}\to\infty$ --- vacuum diffuses infinitely fast (but the damping also vanishes, so the overdamped approximation breaks down at $q\approx 1$). As $q\to 0$, $D_{\rm eff}\to c^2/\gamma_0 = ca$, a finite value --- the minimum diffusivity is set by the Planck-scale constants.

\textbf{Condition for overdamped regime:} The crossover is controlled by the dimensionless ratio:
\begin{equation}
\varepsilon \equiv \frac{\omega}{\gamma_0(1-q)},
\end{equation}
where $\omega$ is a characteristic frequency. The overdamped regime corresponds to $\varepsilon \ll 1$. This replaces the ``$\gamma\to\infty$'' singular perturbation of earlier versions with the physically meaningful small parameter $\varepsilon$.

\subsubsection{A3.7 Wave speed}

For small perturbations in vacuum ($q\approx 1$), the telegraph equation reduces to $\Box(\delta q)=0$, whose characteristic speed is exactly $c$. There is no $q_0$ factor, no $\kappa/\tau$ ratio to calibrate --- the propagation speed is $c$ by construction.

\textbf{GW170817 consistency:} The observed gravitational wave speed $v_{\rm GW} = c(1 \pm 10^{-15})$ is automatically satisfied: in vacuum, $\gamma_0(1-q)=0$, and the dynamics are exactly $\Box q=0$ with speed $c$. The small deviation from $c$ arises only from the nonzero $(1-q)$ in the intergalactic medium; see \S D4.

\textbf{Comparison with constant-damping formulation:} With a constant damping coefficient $\gamma$, the telegraph equation takes the form $\partial_t^2 q + \gamma\partial_t q = (\kappa/\tau)\nabla^2(\ln q)$ with signal speed $c_s = \sqrt{\kappa/(\tau q_0)}$, which requires calibration $\kappa/\tau = c^2 q_0$ to match observations. In the present unified formulation, $c$ appears directly as the propagation speed, the damping is $\gamma_0(1-q)$ rather than a constant $\gamma$, and no ad hoc calibration of the wave speed is needed.

\textbf{Approximation flagged:} The continuum limit neglects the $\mathbf{J}\cdot\nabla q$ cross-term in deriving the closed telegraph equation. This term is $\mathcal{O}(|\nabla q|^2)$ and is negligible for small perturbations. For strong gradients (near $q\to 0$ cores), the discrete treatment is necessary.

% ============================================================
\subsection{A4. Full Expansion of $\nabla^2(\ln q) = \nabla^2 q/q - |\nabla q|^2/q^2$}
% ============================================================

\subsubsection{A4.1 Vector identity derivation}

Let $\psi \equiv \ln q$, so $q = e^\psi$. By the chain rule:
\begin{equation}
\nabla q = \nabla(e^\psi) = e^\psi \nabla\psi = q\,\nabla\psi.
\end{equation}

Therefore:
\begin{equation}
\nabla\psi = \frac{\nabla q}{q}.
\end{equation}

Take the divergence:
\begin{align}
\nabla^2\psi = \nabla\cdot(\nabla\psi) &= \nabla\cdot\left(\frac{\nabla q}{q}\right)\\
&= \frac{\nabla\cdot(\nabla q)}{q} + (\nabla q)\cdot\nabla\left(\frac{1}{q}\right)\\
&= \frac{\nabla^2 q}{q} + (\nabla q)\cdot\left(-\frac{\nabla q}{q^2}\right)\\
&= \frac{\nabla^2 q}{q} - \frac{|\nabla q|^2}{q^2}.
\end{align}

Hence:
\begin{equation}
\boxed{\nabla^2(\ln q) = \frac{\nabla^2 q}{q} - \frac{|\nabla q|^2}{q^2}}.
\label{eq:laplacian_expansion}
\end{equation}

\subsubsection{A4.2 Equivalent forms}

The unified telegraph equation can be written in three equivalent forms:
\begin{align}
\partial_t^2 q + \gamma_0(1-q)\partial_t q &= c^2\,\nabla^2(\ln q) &&\text{(compact form)}, \\
\partial_t^2 q + \gamma_0(1-q)\partial_t q &= c^2\left(\frac{\nabla^2 q}{q} - \frac{|\nabla q|^2}{q^2}\right) &&\text{(expanded form)}, \\
\partial_t^2 q + \gamma_0(1-q)\partial_t q &= c^2\,\nabla\cdot\left(\frac{\nabla q}{q}\right) &&\text{(divergence form)}.
\end{align}

\subsubsection{A4.3 The $1/q$ amplification mechanism}

The expanded form reveals the physical mechanism:
\begin{itemize}
\item The term $\nabla^2 q/q$ is the standard diffusion term, amplified by $1/q$: in low-$q$ regions, even small Laplacians produce large effects.
\item The term $-|\nabla q|^2/q^2$ is a nonlinear self-organization term: it amplifies existing gradients (always non-positive contribution), driving $q$ toward further inhomogeneity.
\end{itemize}

In static equilibrium, the two balance: $\nabla^2 q/q = |\nabla q|^2/q^2$, i.e., $\nabla^2(\ln q)=0$. This is the condition for a self-organized critical state.

\textbf{No approximations in A4.} Equation (\ref{eq:laplacian_expansion}) is an exact vector calculus identity valid wherever $q>0$ and $q$ is twice differentiable.

% ============================================================
\subsection{A5. Spherical Static Solution: From $\nabla^2(\ln 1/q)=0$ to $q(r)=e^{-GM/rc^2}$}
% ============================================================

\subsubsection{A5.1 Static equation in terms of $\psi$}

Define the information potential $\psi \equiv -\ln q = \ln(1/q)$. The static equation $\nabla^2(\ln q)=0$ is equivalent to:
\begin{equation}
\nabla^2\psi = 0.
\label{eq:static_psi}
\end{equation}

This is the Laplace equation for $\psi$.

\subsubsection{A5.2 Spherical symmetry}

Assume $\psi = \psi(r)$, $r=|\mathbf{x}|$. The Laplace operator in spherical coordinates:
\begin{equation}
\nabla^2\psi = \frac{1}{r^2}\frac{d}{dr}\left(r^2\frac{d\psi}{dr}\right) = 0.
\end{equation}

Multiply by $r^2$ and integrate once:
\begin{equation}
r^2\frac{d\psi}{dr} = C_1 = \text{const}.
\label{eq:first_integral}
\end{equation}

\subsubsection{A5.3 Determining $C_1$ via Gauss's law with source}

To introduce a source, we generalize the static equation to include a source term representing locked information (mass):
\begin{equation}
\nabla^2\psi = -\frac{4\pi G}{c^2}\,\rho_m(\mathbf{x}),
\label{eq:poisson_psi}
\end{equation}
where $\rho_m$ is the mass density. The coupling $4\pi G/c^2$ is calibrated to match Newtonian gravity (see \S A6). For a point mass $M$ at the origin, $\rho_m(\mathbf{x}) = M\delta^{(3)}(\mathbf{x})$.

Integrate (\ref{eq:poisson_psi}) over a sphere of radius $r$:
\begin{equation}
\int_{B_r} \nabla^2\psi\,d^3x = -\frac{4\pi G}{c^2}M.
\end{equation}

By the divergence theorem:
\begin{equation}
\oint_{S_r} \nabla\psi\cdot d\mathbf{S} = 4\pi r^2\frac{d\psi}{dr} = -\frac{4\pi G}{c^2}M.
\end{equation}

Therefore:
\begin{equation}
\frac{d\psi}{dr} = -\frac{GM}{c^2 r^2}, \qquad C_1 = -\frac{GM}{c^2}.
\label{eq:C1_value}
\end{equation}

\textbf{Key point:} Equation (\ref{eq:poisson_psi}) is the \emph{definition} of how mass sources the $\psi$-field. The $1/r^2$ force law follows, but the proportionality constant $G$ is calibrated, not derived.

\subsubsection{A5.4 Second integration}

From $d\psi/dr = -GM/(c^2 r^2)$:
\begin{equation}
\psi(r) = \frac{GM}{c^2 r} + C_2.
\end{equation}

\subsubsection{A5.5 Boundary condition: $q(\infty)=1$}

At spatial infinity, the influence of any finite mass vanishes: the field should approach the vacuum state $q=1$. Therefore $\psi(\infty) = -\ln(1) = 0$. This forces $C_2 = 0$:
\begin{equation}
\psi(r) = \frac{GM}{c^2 r}.
\end{equation}

\subsubsection{A5.6 Recovering $q(r)$}

Since $\psi = -\ln q$, we have $q = e^{-\psi}$:
\begin{equation}
\boxed{q(r) = e^{-GM/rc^2}}.
\label{eq:q_spherical}
\end{equation}

\subsubsection{A5.7 Properties of the solution}

\begin{enumerate}[label=(\roman*)]
\item $q(r) \to 1$ as $r\to\infty$ (asymptotic flatness).
\item $q(r) \to 0$ as $r\to 0$ (information capacity exhaustion at the center).
\item $q(r) > 0$ for all $r > 0$ (no sharp event horizon --- $q$ never reaches exactly zero at any finite $r$).
\item The solution is self-similar: $q(\lambda r; \lambda M) = q(r; M)$ for any $\lambda > 0$.
\end{enumerate}

\subsubsection{A5.8 Universality}

For any flux potential $\phi\in\Phi$ satisfying $\phi'(q)<0$ and $\phi(0^+)=+\infty$, the static equation $\nabla^2\phi(q)=0$ yields a $1/r$ asymptotic decay in $d=3$ dimensions. The exponential form is specific to $\phi(q)=-\ln q$; other members of $\Phi$ give different functional forms (e.g., $\phi(q)=1/q$ gives $q(r)=1/(1+GM/rc^2)$). All produce identical qualitative physics: $1/r$ decay, no sharp horizon, approach to $q=0$ as $r\to 0$.

\textbf{Approximation flagged:} The derivation assumes a \textbf{point mass} source. For extended mass distributions, one solves $\nabla^2\psi = -(4\pi G/c^2)\rho_m(\mathbf{x})$ with $\rho_m$ extended, and the solution is the convolution of the point-mass Green's function with $\rho_m$:
\begin{equation}
\psi(\mathbf{x}) = \frac{G}{c^2}\int \frac{\rho_m(\mathbf{x}')}{|\mathbf{x}-\mathbf{x}'|}\,d^3x'.
\end{equation}

% ============================================================
\subsection{A6. Weak-Field Expansion: $q(r)\approx 1-GM/rc^2$ $\to$ Newtonian Limit}
% ============================================================

\subsubsection{A6.1 Taylor expansion}

For $GM/rc^2 \ll 1$ (weak field, large $r$), expand the exponential:
\begin{equation}
q(r) = e^{-GM/rc^2} = 1 - \frac{GM}{rc^2} + \frac{1}{2}\left(\frac{GM}{rc^2}\right)^2 - \cdots
\end{equation}

To first order:
\begin{equation}
\boxed{q(r) \approx 1 - \frac{GM}{rc^2}}.
\label{eq:weak_field}
\end{equation}

\subsubsection{A6.2 Newtonian gravitational potential}

The Newtonian gravitational potential for a point mass is:
\begin{equation}
\Phi_N(r) = -\frac{GM}{r}.
\end{equation}

Comparing with (\ref{eq:weak_field}):
\begin{equation}
q(r) \approx 1 + \frac{\Phi_N(r)}{c^2}.
\end{equation}

Thus, to first order in the weak-field expansion:
\begin{equation}
\boxed{\frac{\Phi_N}{c^2} = q - 1 = -(1-q)}.
\label{eq:newton_correspondence}
\end{equation}

The Newtonian potential (in units of $c^2$) equals \emph{minus the deviation of $q$ from unity}---i.e., minus the occupied fraction $1-q$.

\subsubsection{A6.3 Gravitational acceleration}

The gravitational acceleration is:
\begin{equation}
\mathbf{g} = -\nabla\Phi_N = -c^2\nabla q.
\end{equation}

For the spherical solution: $\mathbf{g} = -c^2\,d(e^{-GM/rc^2})/dr\,\hat{\mathbf{r}} = -(GM/r^2)e^{-GM/rc^2}\,\hat{\mathbf{r}} \approx -(GM/r^2)\,\hat{\mathbf{r}}$ in the weak-field limit. This recovers the inverse-square law.

\subsubsection{A6.4 Poisson equation recovery}

From $\Phi_N = c^2(q-1)$ and the static $q$-field equation $\nabla^2(\ln q)=0$:

Expand $\ln q = \ln(1 + \Phi_N/c^2) \approx \Phi_N/c^2$ for $|\Phi_N|/c^2 \ll 1$. Then in vacuum $\nabla^2(\ln q) \approx \nabla^2(\Phi_N)/c^2 = 0$, so $\nabla^2\Phi_N=0$. With sources (as derived in \S A5.3):
\begin{equation}
\nabla^2\Phi_N = 4\pi G\rho_m,
\end{equation}
which is exactly the Newtonian Poisson equation.

\subsubsection{A6.5 Validity regime}

The weak-field approximation requires:
\begin{equation}
\frac{GM}{rc^2} \ll 1 \quad\Longleftrightarrow\quad r \gg \frac{GM}{c^2} \equiv \frac{r_s}{2},
\end{equation}
where $r_s = 2GM/c^2$ is the Schwarzschild radius. For the Earth, $GM/rc^2 \sim 7\times 10^{-10}$ at the surface. For a neutron star, $GM/Rc^2 \sim 0.21$. For a black hole at the horizon, $GM/r_s c^2 = 1/2$.

\textbf{Approximation flagged:} The identification $\Phi_N = c^2(q-1)$ is valid only to first order in $GM/rc^2$. The next-to-leading-order term differs between potential forms, and neither can be trusted without a full strong-field treatment.

% ============================================================
\subsection{A7. Generalized Lyapunov Functional for Capacity-Dependent Damping}
% ============================================================

\subsubsection{A7.1 Motivation}

For the parabolic $q$-field equation $\partial_t q = \kappa\nabla^2(\ln q)$, the free information functional $\mathcal{F}[q] = \int \Phi(q)\,dx$ with $\Phi(q)=q\ln q - q + 1$ is a Lyapunov functional (proved in the continuum limit; see \S A7 for the Lyapunov construction). For the unified telegraph equation (\ref{eq:a3_telegraph}), we require a Lyapunov functional that accounts for the kinetic energy of the current $\mathbf{J}$ and the capacity-dependent damping $\gamma_0(1-q)$. This section proves that such a functional exists for both the full nonlinear and linearized cases.

\subsubsection{A7.2 Definition of the generalized energy functional}

Define:
\begin{equation}
E[q,\mathbf{J}] \equiv \int_{\mathbb{R}^d} \left[\Phi(q) + \frac{|\mathbf{J}|^2}{2c^2}\right] d^d x,
\label{eq:energy_def}
\end{equation}
where $\Phi(q) = q\ln q - q + 1$ (same as in the parabolic case, satisfying $\Phi'(q)=\ln q$, $\Phi(1)=0$, $\Phi(q)\ge 0$). The second term is the ``kinetic energy'' of the information current. The factor $2c^2$ in the denominator is fixed by the continuum current equation (\ref{eq:a3_continuum_current}).

\subsubsection{A7.3 Time derivative --- full nonlinear case}

Compute $dE/dt$:
\begin{align}
\frac{dE}{dt} &= \int \left[\Phi'(q)\partial_t q + \frac{1}{c^2}\mathbf{J}\cdot\partial_t\mathbf{J}\right] d^d x \nonumber \\
&= \int \left[\ln q\,\nabla\cdot\mathbf{J} + \frac{1}{c^2}\mathbf{J}\cdot\left(c^2\nabla(\ln q) - \gamma_0(1-q)\mathbf{J}\right)\right] d^d x \nonumber \\
&= \int \left[\ln q\,\nabla\cdot\mathbf{J} + \mathbf{J}\cdot\nabla(\ln q) - \frac{\gamma_0(1-q)}{c^2}|\mathbf{J}|^2\right] d^d x.
\end{align}

The first two terms combine into a total divergence:
\begin{equation}
\ln q\,\nabla\cdot\mathbf{J} + \mathbf{J}\cdot\nabla(\ln q) = \nabla\cdot(\ln q\,\mathbf{J}).
\end{equation}

Therefore:
\begin{equation}
\frac{dE}{dt} = \oint_{S_\infty} (\ln q)\,\mathbf{J}\cdot d\mathbf{S} - \frac{\gamma_0}{c^2}\int_{\mathbb{R}^d} (1-q)|\mathbf{J}|^2\,d^d x.
\end{equation}

For an isolated system ($q\to 1$, $\mathbf{J}\to 0$ at infinity), the boundary term vanishes. Hence:
\begin{equation}
\boxed{\frac{dE}{dt} = -\frac{\gamma_0}{c^2}\int_{\mathbb{R}^d} (1-q)|\mathbf{J}|^2\,d^d x \le 0}.
\label{eq:energy_dissipation_nonlinear}
\end{equation}

\textbf{This is a strict Lyapunov functional.} The integrand $(1-q)|\mathbf{J}|^2 \ge 0$ (since $q\le 1$). Dissipation ceases if and only if either (i) $\mathbf{J}=0$ everywhere, or (ii) $q=1$ everywhere (pure vacuum, where there is nothing to dissipate). The second case is physically expected: in vacuum, no capacity is occupied, so no current encounters resistance.

\subsubsection{A7.4 Time derivative --- linearized case ($q\approx 1$)}

Near vacuum, $\gamma_0(1-q) \approx 0$, and the damping term vanishes to leading order. The continuum current equation becomes:
\begin{equation}
\partial_t\mathbf{J} \approx c^2\nabla(\delta q),
\end{equation}
where we used $\nabla(\ln q) \approx \nabla(\delta q)$. The telegraph equation reduces to $\Box(\delta q)=0$.

Define the linearized energy:
\begin{equation}
E_{\rm lin}[\delta q, \mathbf{J}] \equiv \frac{1}{2}\int \left[(\partial_t\delta q)^2 + c^2|\nabla(\delta q)|^2 + \frac{|\mathbf{J}|^2}{c^2}\right] d^d x.
\end{equation}

Using $\partial_t(\delta q) = \nabla\cdot\mathbf{J}$ and $\partial_t\mathbf{J} = c^2\nabla(\delta q)$:
\begin{align}
\frac{dE_{\rm lin}}{dt} &= \int \left[\partial_t(\delta q)\,\partial_t^2(\delta q) + c^2\nabla(\delta q)\cdot\nabla\partial_t(\delta q) + \frac{1}{c^2}\mathbf{J}\cdot\partial_t\mathbf{J}\right] d^d x \nonumber \\
&= \int \left[\nabla\cdot\mathbf{J}\,\partial_t^2 q + c^2\nabla q\cdot\nabla(\nabla\cdot\mathbf{J}) + \mathbf{J}\cdot\nabla q\right] d^d x \nonumber \\
&= \oint (\cdots)\cdot d\mathbf{S} = 0,
\end{align}
where all terms combine into boundary integrals that vanish for isolated systems. Energy is conserved in the linearized (vacuum) limit, as expected for a wave equation without dissipation.

\subsubsection{A7.5 Physical interpretation}

\begin{enumerate}[label=(\roman*)]
\item $E[q,\mathbf{J}]$ is a monotonically non-increasing functional. The system evolves irreversibly toward configurations where $\mathbf{J}=0$ wherever $q<1$.
\item Dissipation is proportional to $(1-q)$, the occupied fraction. Where cells are occupied, information currents are obstructed; where cells are vacant ($q\approx 1$), currents flow freely.
\item In vacuum ($q=1$), $dE/dt = 0$ exactly --- energy is conserved. This is the regime where Lorentz invariance and wave dynamics emerge.
\item This is an $H$-theorem for the $q$-field with capacity-dependent damping. The arrow of time is built into the dynamics through the $(1-q)$ factor.
\item The result generalizes the constant-$\gamma$ Cattaneo $H$-theorem to the physically motivated $\gamma(q)=\gamma_0(1-q)$ damping (see \S A7 for the Lyapunov construction).
\end{enumerate}

\subsubsection{A7.6 Generalization to arbitrary damping profile}

If $\gamma(q) = \gamma_0\,h(1-q)$ for any non-negative function $h$ with $h(0)=0$, the same proof structure holds:
\begin{equation}
\frac{dE}{dt} = -\frac{\gamma_0}{c^2}\int h(1-q)|\mathbf{J}|^2\,d^d x \le 0.
\end{equation}
The Lyapunov property is universal across the class of capacity-dependent damping functions. The linear form $h(x)=x$ is the simplest member; the essential physics ($dE/dt\le 0$, vanishing dissipation in vacuum) is independent of this choice.

\textbf{No approximations in A7.} The dissipation theorems are exact consequences of the unified current equation (\ref{eq:a3_current}) and continuity (\ref{eq:continuity_continuum}), assuming sufficient regularity for integration by parts.

%%%%%%%%%%%%%%%%%%%%%%%%%%%%%%%%%%%%%%%%%%%%%%%%%%%%%%%%%%%%%%%%%%%%%
\section{Mathematical Rigor}
%%%%%%%%%%%%%%%%%%%%%%%%%%%%%%%%%%%%%%%%%%%%%%%%%%%%%%%%%%%%%%%%%%%%%

% ============================================================
\subsection{B1. Well-Posedness Proof for the Coupled $(q,\mathbf{J})$ System}
% ============================================================

\subsubsection{B1.1 The system}

We consider the unified hyperbolic system on $\mathbb{R}^d$:
\begin{align}
\partial_t q &= \nabla\cdot\mathbf{J}, \label{eq:sys1_v3}\\
\partial_t\mathbf{J} &= c^2\nabla(\ln q) - \gamma_0(1-q)\mathbf{J}. \label{eq:sys2_v3}
\end{align}
with initial conditions $q(\mathbf{x},0)=q_0(\mathbf{x})$, $\mathbf{J}(\mathbf{x},0)=\mathbf{J}_0(\mathbf{x})$.

\subsubsection{B1.2 Local well-posedness}

\textbf{Theorem B1.} Let $q_0 \in H^s(\mathbb{R}^d)$ with $q_0(\mathbf{x}) \ge q_{\min} > 0$ for all $\mathbf{x}$, and $\mathbf{J}_0 \in H^{s-1}(\mathbb{R}^d)$ with $s > d/2+1$. Then there exists $T>0$ such that the system (\ref{eq:sys1_v3})-(\ref{eq:sys2_v3}) has a unique solution with:
\begin{equation}
q \in C([0,T]; H^s(\mathbb{R}^d)) \cap C^1([0,T]; H^{s-2}(\mathbb{R}^d)), \quad \mathbf{J} \in C([0,T]; H^{s-1}(\mathbb{R}^d)).
\end{equation}

\textbf{Proof sketch.} Rewrite as a first-order symmetric hyperbolic system in $(q, \mathbf{J})$. The nonlinearity $\nabla(\ln q) = \nabla q/q$ is smooth for $q \ge q_{\min} > 0$. The damping term $-\gamma_0(1-q)\mathbf{J}$ is a smooth, bounded perturbation. Standard theory (Kato, Majda) guarantees local existence and uniqueness.

The strict positivity condition $q_0 \ge q_{\min} > 0$ is essential: as $q\to 0$, the nonlinearity loses Lipschitz continuity.

\subsubsection{B1.3 Global existence}

\textbf{Theorem B1 (continued).} If $q_0 \ge \delta > 0$ uniformly, the solution exists for all $t > 0$ and satisfies:
\begin{equation}
E(t) + \frac{\gamma_0}{c^2}\int_0^t \int (1-q)|\mathbf{J}|^2\,d^d x\,dt' = E(0).
\end{equation}

\textbf{Proof sketch.} The energy $E(t)$ (defined in \S A7.2) is non-increasing. Since $\Phi(q) \ge 0$ and $|\mathbf{J}|^2 \ge 0$, both $q$ and $\mathbf{J}$ remain bounded in the energy norm. A continuation argument extends the local solution to global.

\subsubsection{B1.4 The $q\to 0$ problem}

When $\inf q_0 = 0$, the hyperbolic system is not locally well-posed in Sobolev spaces. The formation of $q=0$ singularities corresponds to ``classical cores''---regions where the continuum description breaks down.

\textbf{Flag:} Rigorous well-posedness is proved only for $q \ge q_{\min} > 0$. The formation and dynamics of $q=0$ cores require the discrete theory.

% ============================================================
\subsection{B2. Spectral Analysis: Eigenvalues, Stability, Dispersion Relation}
% ============================================================

\subsubsection{B2.1 Linearization around $q=1$}

Linearize the unified telegraph equation (\ref{eq:a3_telegraph}) around vacuum $q=1$. Set $q = 1 + \delta q$ with $|\delta q|\ll 1$. Then $\ln q \approx \delta q$ and $\gamma_0(1-q)\partial_t q \approx \mathcal{O}(\delta q^2)$. To linear order:
\begin{equation}
\partial_t^2(\delta q) = c^2\nabla^2(\delta q).
\label{eq:linearized_b2}
\end{equation}

This is the wave equation, not the heat equation. The linearization is exact (damping is second order).

\subsubsection{B2.2 Fourier modes and wave dispersion}

Seek solutions $\delta q(\mathbf{x},t) = \hat{q}_{\mathbf{k}}\,e^{i(\mathbf{k}\cdot\mathbf{x} - \omega t)}$. Substituting into (\ref{eq:linearized_b2}):
\begin{equation}
-\omega^2 = -c^2|\mathbf{k}|^2 \quad\Longrightarrow\quad \boxed{\omega(\mathbf{k}) = \pm c|\mathbf{k}|}.
\label{eq:wave_dispersion}
\end{equation}

All modes propagate at speed $c$ without damping. This is the vacuum dispersion of the $q$-field.

\subsubsection{B2.3 Full nonlinear dispersion (local approximation)}

For the full telegraph equation (\ref{eq:a3_telegraph}), treat $q$ as locally constant for the purpose of a WKB-type dispersion analysis. Seek $q = q_0 + \hat{q}\,e^{i(\mathbf{k}\cdot\mathbf{x} - \omega t)}$ with $q_0$ the local background. Then $\nabla^2(\ln q) \approx -|\mathbf{k}|^2\hat{q}/q_0$ and $\partial_t q \approx -i\omega\hat{q}$, $\partial_t^2 q \approx -\omega^2\hat{q}$:
\begin{equation}
-\omega^2 - i\gamma_0(1-q_0)\omega + \frac{c^2}{q_0}|\mathbf{k}|^2 = 0.
\end{equation}

Solving for complex $\omega = \omega_R + i\omega_I$:
\begin{equation}
\omega = -\frac{i}{2}\gamma_0(1-q_0) \pm \sqrt{\frac{c^2}{q_0}|\mathbf{k}|^2 - \frac{\gamma_0^2(1-q_0)^2}{4}}.
\label{eq:nonlinear_dispersion}
\end{equation}

\textbf{Dispersion regimes:}
\begin{enumerate}[label=(\roman*)]
\item \textbf{Wave regime} ($|\mathbf{k}| > k_{\rm cross}$): The square root is real. $\omega_R = \pm\sqrt{c^2|\mathbf{k}|^2/q_0 - \gamma_0^2(1-q_0)^2/4}$, $\omega_I = -\gamma_0(1-q_0)/2$. Damped propagating waves.
\item \textbf{Diffusive regime} ($|\mathbf{k}| < k_{\rm cross}$): The square root is imaginary. Both roots are purely imaginary (non-propagating, purely decaying modes). Overdamped diffusion.
\item \textbf{Crossover:} $k_{\rm cross}(q_0) = \gamma_0(1-q_0)\sqrt{q_0}/(2c)$.
\end{enumerate}

For $q_0\to 1$ (vacuum): $k_{\rm cross} \to 0$. \textbf{All modes are in the wave regime near vacuum.} For $q_0\to 0$ (near saturation): $k_{\rm cross} \to 0$ as well (since $(1-q_0)\sqrt{q_0} \to 0$). The maximum crossover occurs at $q_0 = 1/3$, giving $k_{\rm cross}^{\max} \approx 0.385 \times \gamma_0/(2c) \approx 1.19\times 10^{34}$~m$^{-1}$ --- the crossover wavenumber is at most Planckian.

\subsubsection{B2.4 Eigenvalue problem for the static equation}

In a bounded domain $\Omega$ with boundary conditions $q|_{\partial\Omega}=1$ (or Neumann $\partial_n q|_{\partial\Omega}=0$), the static equation $\nabla^2(\ln q)=0$ with $\ln q\equiv\psi$ is the Dirichlet (or Neumann) eigenvalue problem for the Laplacian. The eigenvalues $\lambda_n$ of $-\nabla^2$ satisfy $0=\lambda_0<\lambda_1\le\lambda_2\le\cdots$. The static solution exists uniquely for Dirichlet conditions; for Neumann conditions, it exists up to an additive constant in $\psi$, fixed by total $q$ conservation.

\textbf{No approximations in B2.} The vacuum dispersion is exact. The nonlinear dispersion is a local (WKB) approximation valid when $q$ varies slowly compared to the wavelength.

% ============================================================
\subsection{B3. Overdamped Limit as Singular Perturbation ($\varepsilon \equiv \omega/[\gamma_0(1-q)] \ll 1$)}
% ============================================================

\subsubsection{B3.1 The genuine small parameter}

In earlier versions of the SM, the overdamped limit was taken as ``$\gamma\to\infty$ with $c_s^2/\gamma$ fixed.'' This was formally correct but physically opaque: $\gamma_0 \sim 1/t_P \sim 10^{43}$~s$^{-1}$ is already enormous. The physically meaningful small parameter is not the absolute value of $\gamma_0$, but the dimensionless ratio:
\begin{equation}
\boxed{\varepsilon \equiv \frac{\omega}{\gamma_0(1-q)}},
\label{eq:epsilon_def}
\end{equation}
where $\omega$ is a characteristic frequency of the dynamics. The overdamped regime is $\varepsilon \ll 1$.

\subsubsection{B3.2 Scaling analysis}

From the telegraph equation (\ref{eq:a3_telegraph}), the relative magnitude of inertial to damping terms is:
\begin{equation}
\frac{|\partial_t^2 q|}{|\gamma_0(1-q)\partial_t q|} \sim \frac{\omega^2}{\gamma_0(1-q)\omega} = \frac{\omega}{\gamma_0(1-q)} = \varepsilon.
\end{equation}

When $\varepsilon \ll 1$, the inertial term is negligible:
\begin{equation}
\gamma_0(1-q)\partial_t q \approx c^2\nabla^2(\ln q),
\end{equation}
which is the overdamped (parabolic-type) equation.

\subsubsection{B3.3 When is $\varepsilon \ll 1$?}

\begin{itemize}
\item \textbf{Vacuum ($q\approx 1$):} $\gamma_0(1-q) \approx 0$, so $\varepsilon \to \infty$ regardless of $\omega$. The system is \textbf{always underdamped} near vacuum. The wave equation $\Box q = 0$ is the correct description.
\item \textbf{Laboratory ($q\approx 1-10^{-27}$):} $\gamma_0(1-q) \sim 10^{16}$~s$^{-1}$. For $\omega \sim 1$~Hz (typical lab timescale), $\varepsilon \sim 10^{-16}$. \textbf{Deeply overdamped.} The parabolic approximation is excellent.
\item \textbf{Solar surface ($q\approx 1-2\times 10^{-6}$):} $\gamma_0(1-q) \sim 4\times 10^{37}$~s$^{-1}$. For $\omega \sim 10^{-3}$~Hz (orbital), $\varepsilon \sim 2.5\times 10^{-41}$. \textbf{Extraordinarily overdamped.}
\item \textbf{Near black hole ($q\sim 0.6$, $\gamma_0(1-q)\sim 7\times 10^{42}$~s$^{-1}$):} For $\omega \sim c/r_s \sim 10^{4}$~Hz (dynamical), $\varepsilon \sim 1.4\times 10^{-39}$. Still deeply overdamped.
\end{itemize}

\subsubsection{B3.4 Singular perturbation structure}

The overdamped limit is a singular perturbation: the order of the PDE drops from 2 (hyperbolic) to 1 (parabolic-type). The second initial condition (initial current $\mathbf{J}_0$) is lost in a boundary layer of thickness $\mathcal{O}(\varepsilon/\omega) = \mathcal{O}(1/[\gamma_0(1-q)])$ at $t=0$.

Formally, the solution has the asymptotic expansion:
\begin{equation}
q(\mathbf{x},t;\varepsilon) = q^{(0)}(\mathbf{x},t) + e^{-\gamma_0(1-q)t}q^{(1)}(\mathbf{x},t) + \mathcal{O}(\varepsilon),
\end{equation}
where $q^{(0)}$ solves the overdamped equation and the initial layer $q^{(1)}$ decays exponentially on the fast timescale $1/[\gamma_0(1-q)]$.

\textbf{Physical interpretation:} On timescales $t \gg 1/[\gamma_0(1-q)]$, the system ``forgets'' the initial current configuration and follows the gradient-driven overdamped dynamics. For laboratory and astrophysical timescales, the overdamped limit is an extraordinarily good approximation --- \emph{except} in vacuum ($q\approx 1$), where the overdamped approximation fails and the full wave dynamics are mandatory.

The singular perturbation is expressed in terms of the physically meaningful $\varepsilon = \omega/[\gamma_0(1-q)]$ rather than the formal limit $\gamma\to\infty$, which obscured the fact that near vacuum the system is \emph{not} overdamped at any frequency.

% ============================================================
\subsection{B4. Continuum Limit of Graph Laplacian on Locally Finite Graphs}
% ============================================================

\subsubsection{B4.1 Discrete setup}

Consider an infinite regular grid $\Lambda \subset \mathbb{Z}^d$ with spacing $a$. The graph Laplacian acting on a function $u_i = u(\mathbf{x}_i)$ is:
\begin{equation}
(\mathcal{L}_a u)_i \equiv \frac{1}{a^2}\sum_{j \sim i} (u_j - u_i),
\end{equation}
where $j\sim i$ denotes nearest neighbors.

\subsubsection{B4.2 Taylor expansion}

For a smooth function $u(\mathbf{x})$, Taylor-expand around $\mathbf{x}_i$:
\begin{equation}
u(\mathbf{x}_i \pm a\hat{\mathbf{e}}_\alpha) = u(\mathbf{x}_i) \pm a\partial_\alpha u(\mathbf{x}_i) + \frac{a^2}{2}\partial_\alpha^2 u(\mathbf{x}_i) \pm \frac{a^3}{6}\partial_\alpha^3 u(\mathbf{x}_i) + \mathcal{O}(a^4).
\end{equation}

Summing over $\alpha = 1,\ldots,d$, with two neighbors per direction:
\begin{align}
(\mathcal{L}_a u)_i &= \frac{1}{a^2}\sum_{\alpha=1}^d \Big[(u(\mathbf{x}_i + a\hat{\mathbf{e}}_\alpha) - u(\mathbf{x}_i)) + (u(\mathbf{x}_i - a\hat{\mathbf{e}}_\alpha) - u(\mathbf{x}_i))\Big]\\
&= \frac{1}{a^2}\sum_{\alpha=1}^d \left[a^2\partial_\alpha^2 u + \mathcal{O}(a^4)\right]\\
&= \nabla^2 u(\mathbf{x}_i) + \mathcal{O}(a^2).
\label{eq:graph_laplacian_limit}
\end{align}

The $\mathcal{O}(a^3)$ terms cancel by symmetry. The leading error is $\mathcal{O}(a^2)$.

\subsubsection{B4.3 Locally finite graphs}

A graph is \textbf{locally finite} if every vertex has finite degree. For the $q$-field model, the degree is $2d$ (regular grid). The result (\ref{eq:graph_laplacian_limit}) generalizes to any locally finite graph with a periodic structure, provided the graph is ``sufficiently uniform'' (bounded geometry, polynomial volume growth).

\subsubsection{B4.4 The $q$-field PDE from the discrete dynamics}

The discrete $q$-field dynamics follows from the unified current equation. In the static limit, $\partial_t J_{ij}=0$ gives $J_{ij} \propto (\ln q_i - \ln q_j)$. The no-net-current condition $\sum_j J_{ij}=0$ yields $\sum_j (\ln q_i - \ln q_j) = 0$, which in the continuum becomes $\nabla^2(\ln q)=0$.

\textbf{Approximation flagged:} The convergence $\mathcal{L}_a \to \nabla^2$ is pointwise in $C^4$ functions. Near $q=0$ cores, the higher-order terms may be significant.

% ============================================================
\subsection{B5. Discrete-to-Continuum Error Estimates ($N\sim 10^{184}$ Cells)}
% ============================================================

\subsubsection{B5.1 Cell counting --- corrected from $10^{90}$ to $10^{184}$}

\textbf{CORRECTION:} An earlier estimate of $N\sim 10^{90}$ used a nuclear-scale cell size $\sim 10^{-15}$~m, which is the scale of quantum field theory degrees of freedom, not the DGF cell scale. The DGF cell spacing is identified with the Planck length $\ell_P$ (see \S F3). The correct estimate is:

The observable universe has volume $V \sim (4\pi/3)(c/H_0)^3 \sim 10^{80}$ m$^3$. If the fundamental cell has Planck volume $\ell_P^3 \sim (1.6\times 10^{-35}~\text{m})^3 \sim 4\times 10^{-105}$ m$^3$, then:
\begin{equation}
\boxed{N = \frac{V}{\ell_P^3} \sim \frac{10^{80}}{4\times 10^{-105}} \sim 2.5\times 10^{184}}.
\label{eq:N_estimate}
\end{equation}

This is a fiducial estimate. The actual number is not derived from Definition~C and may differ by many orders of magnitude. However, as long as $N \gg 10^{20}$, the continuum approximation is effectively exact for all macroscopic purposes.

\subsubsection{B5.2 Discretization error scaling}

For a smooth function $u$ on a grid of $N$ cells (spacing $a \sim N^{-1/d}$), the discrete Laplacian approximates the continuum Laplacian with error (from \S B4):
\begin{equation}
\|(\mathcal{L}_a - \nabla^2)u\|_\infty \le C a^2 \|\nabla^4 u\|_\infty \sim \mathcal{O}(N^{-2/d}).
\end{equation}

In $d=3$: error $\sim \mathcal{O}(N^{-2/3})$. For $N\sim 10^{184}$: error $\sim 10^{-123}$. This is the pointwise error per cell.

\subsubsection{B5.3 Global error estimate}

For the static spherical solution $q(r) = e^{-GM/rc^2}$, the fourth derivative scales as $\nabla^4 q \sim GM/(c^2 r^5)$. The error is largest at small $r$. The continuum approximation is reliable when:
\begin{equation}
a^2 \frac{GM}{c^2 r^5} \ll 1 \quad\Longleftrightarrow\quad r \gg (r_s a^2)^{1/5}.
\end{equation}

For a solar-mass black hole ($r_s \sim 3\times 10^3$~m) and Planck-scale cells ($a \sim 10^{-35}$~m): $(r_s a^2)^{1/5} \sim (3\times 10^3 \times 10^{-70})^{1/5} \sim 10^{-13}$~m, far below any physically accessible scale.

\subsubsection{B5.4 Finite-size effects}

For a finite system of $N$ cells, the total number of microstates is $2^N$. Fluctuations in macroscopic observables (like $q$) scale as $1/\sqrt{N}$. For $N\sim 10^{184}$:
\begin{equation}
\frac{\Delta q}{q} \sim \frac{1}{\sqrt{N}} \sim 10^{-92}.
\end{equation}

Finite-size fluctuations are astronomically suppressed.

\subsubsection{B5.5 Important caveat}

\textbf{The error estimate refers to the finite-$N$ correction to the continuum PDE, not to the absolute accuracy of the DGF framework.} If the DGF framework is an approximation to a deeper theory, this error estimate does not capture that mismatch. The $\mathcal{O}(10^{-123})$ estimate says: \emph{if} the discrete DGF model on $N=10^{184}$ cells is the exact microscopic theory, \emph{then} the continuum PDE approximates it with error $\mathcal{O}(10^{-123})$.

\textbf{Flag:} The cell count $N\sim 10^{184}$ is a fiducial estimate using the Planck scale. The cell spacing $a$ is identified with $\ell_P$ for dimensional consistency (see \S F3), but this identification is an input, not a derivation.

%%%%%%%%%%%%%%%%%%%%%%%%%%%%%%%%%%%%%%%%%%%%%%%%%%%%%%%%%%%%%%%%%%%%%
\section{Physical Precedents}
%%%%%%%%%%%%%%%%%%%%%%%%%%%%%%%%%%%%%%%%%%%%%%%%%%%%%%%%%%%%%%%%%%%%%

The mathematical architecture of the $q$-field---current inertia, finite-speed wave propagation, capacity-dependent damping---has precedents across multiple subfields. These are not analogies that prove the theory's physics; they are existence demonstrations showing that the mathematical structure is sound. Each is discussed below.

% ============================================================
\subsection{C1. Maxwell-Cattaneo Heat Conduction (Second Sound in NaF, He-4)}
% ============================================================

\subsubsection{C1.1 Historical context}

Fourier's law of heat conduction $\mathbf{J}_Q = -k\nabla T$ leads to the parabolic heat equation, which predicts infinite propagation speed. Cattaneo (1948) and Vernotte (1958) independently proposed:
\begin{equation}
\tau\,\partial_t\mathbf{J}_Q + \mathbf{J}_Q = -k\nabla T,
\end{equation}
which leads to the telegraph equation for temperature and predicts finite propagation speed $v = \sqrt{\alpha/\tau}$.

\subsubsection{C1.2 Experimental confirmation}

The wavelike propagation of heat (``second sound'') has been observed in solid NaF (Jackson \& Walker, 1971), liquid He-4 (Peshkov, 1944), and dielectric crystals (Narayanamurti \& Dynes, 1972).

\subsubsection{C1.3 Structural analogy}

The Cattaneo heat equation and the unified $q$-field current equation (\ref{eq:a3_current}) share the identical mathematical structure: a gradient-driven current with finite relaxation, producing a telegraph equation. The physics differs (thermal energy vs.\ information), but the mathematics is isomorphic. This provides an ``existence proof'' that flux with finite relaxation time is physically realizable.

\subsubsection{C1.4 Key difference from the unified framework}

In heat conduction, the Cattaneo relaxation time $\tau$ is a material constant (temperature-dependent but externally specified). In the $q$-field framework, the effective damping $\gamma_0(1-q)$ is \emph{motivated by Definition~C}: occupied cells block new causal events, so damping is proportional to the occupied fraction $(1-q)$. The linear form $\gamma(q)=\gamma_0(1-q)$ is the simplest ansatz consistent with this constraint (flagged as a modeling choice in the main paper, Sec.~IV). The essential property $\gamma\to0$ as $q\to1$ follows from the constraint regardless of the specific functional form.

\textbf{Reference:} D.D.\ Joseph \& L.\ Preziosi, Rev.\ Mod.\ Phys.\ \textbf{61}, 41 (1989); \textbf{62}, 375 (1990).

% ============================================================
\subsection{C2. London Equations (Superconductivity --- Current as Dynamical Variable)}
% ============================================================

The London brothers (1935) elevated the superconducting current $\mathbf{J}_s$ from a constitutive relation to a dynamical field:
\begin{equation}
\frac{\partial}{\partial t}(\Lambda\mathbf{J}_s) = \mathbf{E}, \qquad \Lambda = \frac{m}{n_s e^2}.
\end{equation}

The $q$-field framework makes the identical conceptual move: the current $\mathbf{J}$ is elevated from a constitutive relation to a dynamical field obeying its own equation of motion (\ref{eq:a3_continuum_current}). The crucial difference is that in the London theory, the ``inertia'' parameter $\Lambda$ is a material property; in this framework, the corresponding inertia is fixed by the fundamental speed $c$ and cell spacing $a$, with the damping $\gamma_0(1-q)$ following from Definition~C.

\textbf{Reference:} F.\ London \& H.\ London, Proc.\ R.\ Soc.\ A \textbf{149}, 71 (1935).

% ============================================================
\subsection{C3. Maxwell's Displacement Current ($\partial_t\mathbf{E}$ Term $\to$ Finite Propagation Speed)}
% ============================================================

Maxwell's addition of $\partial_t\mathbf{E}$ to Amp\`ere's law converted electromagnetic equations from parabolic/elliptic to hyperbolic, enabling wave solutions with finite speed $c$. The unified current equation makes the same structural move: the $\partial_t J_{ij}$ term (current inertia) is what enables finite-speed propagation. Without it, the flux would be instantaneously determined by the gradient, and the dynamics would be parabolic.

\textbf{Reference:} J.C.\ Maxwell, Phil.\ Trans.\ R.\ Soc.\ \textbf{155}, 459 (1865).

% ============================================================
\subsection{C4. Heaviside Transmission Line (Telegraph Equation Origin)}
% ============================================================

Heaviside (1876) derived the telegrapher's equation for voltage on a transmission line:
\begin{equation}
LC\,\partial_t^2 V + (LG+RC)\,\partial_t V + RG\,V = \partial_x^2 V.
\end{equation}

For a lossless line ($R=G=0$), this is the wave equation. With damping, it is the telegrapher's equation. The telegraph equation (\ref{eq:a3_telegraph}) has the same structure, with the crucial difference that the damping coefficient $\gamma_0(1-q)$ depends on the field $q$ itself --- a nonlinear transmission line analogy.

\textbf{Reference:} O.\ Heaviside, Phil.\ Mag.\ \textbf{2}, 135 (1876).

% ============================================================
\subsection{C5. Persistent Random Walks (Microscopic Foundation)}
% ============================================================

A persistent random walk (F\"urth, 1917; Taylor, 1921; Goldstein, 1951) introduces correlation between successive steps. The resulting kinetic equation yields, for the spatial density, a telegraph equation with finite propagation speed. The DGF information current corresponds to the net drift of information carriers on the graph. The finite inertia of the current ($\partial_t J_{ij}$ term) reflects the microscopic persistence of information movement.

\textbf{Reference:} G.H.\ Weiss, \emph{Aspects and Applications of the Random Walk} (North-Holland, 1994).

% ============================================================
\subsection{C6. Caldeira-Leggett Model (Wave $\to$ Diffusion Crossover from Bath Coupling)}
% ============================================================

Caldeira and Leggett (1983) studied a quantum particle coupled to a bath of harmonic oscillators. In the high-friction limit, the particle dynamics reduce to diffusion; at finite friction, the dynamics are hyperbolic (Kramers equation). The $q$-field exhibits the identical crossover: the ``bath'' is the ensemble of all other cells, the ``friction'' is $\gamma_0(1-q)$, and the crossover from wave to diffusion is controlled by $(1-q)$ --- precisely the capacity-dependent damping.

\textbf{Reference:} A.O.\ Caldeira \& A.J.\ Leggett, Ann.\ Phys.\ \textbf{149}, 374 (1983).

%%%%%%%%%%%%%%%%%%%%%%%%%%%%%%%%%%%%%%%%%%%%%%%%%%%%%%%%%%%%%%%%%%%%%
\section{Gravitational Wave Details}
%%%%%%%%%%%%%%%%%%%%%%%%%%%%%%%%%%%%%%%%%%%%%%%%%%%%%%%%%%%%%%%%%%%%%

The unified telegraph equation (\ref{eq:a3_telegraph}) is hyperbolic with finite propagation speed. This section works out the observable consequences. There is no separate ``hyperbolic extension'' to postulate here---the telegraph equation is fundamental. What follows is a systematic unpacking of its predictions.

% ============================================================
\subsection{D1. Derivation of Source Term $S(\mathbf{x},t)$ from Binary $q$-Field Quadrupole}
% ============================================================

\subsubsection{D1.1 Wave equation with source}

The static equation with source is $\nabla^2\psi = -(4\pi G/c^2)\rho_m$. The minimal Lorentz-invariant generalization is to replace $\nabla^2$ with the d'Alembertian:
\begin{equation}
\Box\psi \equiv \left(-\frac{1}{c^2}\partial_t^2 + \nabla^2\right)\psi = -\frac{4\pi G}{c^2}\,\rho_m(\mathbf{x},t),
\label{eq:scalar_wave}
\end{equation}
where $\psi = -\ln q$. Note that this is the inhomogeneous wave equation, which follows from perturbing the telegraph equation around $q=1$ with a source: $\Box(\delta q) = (4\pi G/c^2)\rho_m$ to linear order.

\subsubsection{D1.2 Quadrupole formula}

For a non-relativistic source ($v \ll c$), the leading radiating term is the mass quadrupole:
\begin{equation}
\boxed{\psi(\mathbf{x},t) \approx \frac{G}{2c^4 r}\,\hat{n}^i\hat{n}^j\,\ddot{I}_{ij}(t - r/c)},
\label{eq:quadrupole}
\end{equation}
where $I_{ij}(t) \equiv \int \rho_m(\mathbf{x},t)\,x^i x^j\,d^3x$ is the mass quadrupole moment tensor. The monopole term is static ($\sim GM/c^2 r$) and does not radiate. The dipole term vanishes by momentum conservation.

\subsubsection{D1.3 Binary source}

For a binary system of masses $m_1, m_2$ in circular orbit with separation $R$ and orbital frequency $\Omega = \sqrt{G(m_1+m_2)/R^3}$, the quadrupole moment oscillates at frequency $2\Omega$:
\begin{equation}
\ddot{I}_{ij} \propto \mu R^2\Omega^2\,e^{2i\Omega(t - r/c)},
\end{equation}
where $\mu = m_1 m_2/(m_1+m_2)$ is the reduced mass. The $\psi$-field radiation has frequency $f = \Omega/\pi$.

\textbf{Key difference from GR:} In GR, gravitational waves are tensor perturbations $h_{ij}^{\rm TT}$ with two polarization states ($+$, $\times$). In DGF, the radiation is in the scalar field $\psi = -\ln q$. The quadrupole formula for the scalar amplitude differs from GR's tensor quadrupole formula in the numerical prefactor and polarization content.

\textbf{Approximation flagged:} This derivation assumes (i) the wave equation (\ref{eq:scalar_wave}) as the relativistic extension, (ii) minimal coupling to matter via $\rho_m$, and (iii) the quadrupole formula for non-relativistic sources. The coupling to matter requires further development (see OP-3 in \S E3).

% ============================================================
\subsection{D2. Wave-Diffusion Crossover --- Full Dispersion Relation}
% ============================================================

\subsubsection{D2.1 Dispersion relation for the unified telegraph equation}

For plane-wave perturbations $q = q_0 + \hat{q}\,e^{i(\mathbf{k}\cdot\mathbf{x} - \omega t)}$ propagating on a background $q_0$, the unified telegraph equation (\ref{eq:a3_telegraph}) with $\nabla^2(\ln q) \approx -|\mathbf{k}|^2\hat{q}/q_0$ yields:
\begin{equation}
-\omega^2 - i\gamma_0(1-q_0)\omega + \frac{c^2}{q_0}|\mathbf{k}|^2 = 0.
\label{eq:full_dispersion}
\end{equation}

Solving for complex $\omega = \omega_R + i\omega_I$:
\begin{equation}
\omega = -\frac{i\gamma_0(1-q_0)}{2} \pm \sqrt{\frac{c^2}{q_0}|\mathbf{k}|^2 - \frac{\gamma_0^2(1-q_0)^2}{4}}.
\label{eq:omega_complex}
\end{equation}

\subsubsection{D2.2 Crossover wavenumber}

The transition from propagating (real $\omega_R$) to purely decaying (imaginary $\omega_R$) modes occurs when the discriminant vanishes:
\begin{equation}
\boxed{k_{\rm cross}(q_0) = \frac{\gamma_0(1-q_0)}{2c}\sqrt{q_0}}.
\label{eq:kcross_q0}
\end{equation}

For $|\mathbf{k}| > k_{\rm cross}$: damped waves with phase velocity $v_\phi \to c/\sqrt{q_0}$ at high $k$. For $|\mathbf{k}| < k_{\rm cross}$: overdamped, purely decaying modes.

\subsubsection{D2.3 Phase and group velocities}

In the wave regime ($|\mathbf{k}| \gg k_{\rm cross}$):
\begin{align}
\omega_R &\approx \pm \frac{c}{\sqrt{q_0}}|\mathbf{k}|, \\
\omega_I &\approx -\frac{\gamma_0(1-q_0)}{2}.
\end{align}

Phase velocity: $v_\phi = \omega_R/|\mathbf{k}| \approx c/\sqrt{q_0}$.
Group velocity: $v_g = d\omega_R/d|\mathbf{k}| \approx c/\sqrt{q_0}$.

For $q_0 \to 1$ (vacuum): $v_\phi \to c$, $v_g \to c$. Both equal $c$ exactly in the linearized limit.

In the diffusive regime ($|\mathbf{k}| \ll k_{\rm cross}$), expand the square root:
\begin{align}
\omega_+ &\approx -i\frac{c^2|\mathbf{k}|^2}{\gamma_0 q_0(1-q_0)}, \\
\omega_- &\approx -i\gamma_0(1-q_0) + i\frac{c^2|\mathbf{k}|^2}{\gamma_0 q_0(1-q_0)}.
\end{align}

The $\omega_+$ mode is the diffusive mode (purely imaginary, long-lived at small $k$). The $\omega_-$ mode is the rapidly decaying current-relaxation mode.

\subsubsection{D2.4 Vacuum limit}

As $q_0 \to 1$, $k_{\rm cross} \to 0$, and all finite-$k$ modes enter the wave regime. In vacuum, there is no diffusive regime --- all perturbations propagate as waves. This is the mathematical origin of Lorentz emergence: the capacity-dependent damping $\gamma_0(1-q_0)$ vanishes, leaving the pure wave equation.

% ============================================================
\subsection{D3. Quantitative $k_{\rm cross}$ Values Across Astrophysical Environments}
% ============================================================

Using $\gamma_0 = c/\ell_P \approx 1.85\times 10^{43}$~s$^{-1}$:

\begin{table}[h]
\caption{Crossover wavenumber and wavelength for different environments.}
\label{tab:kcross}
\small\setlength{\tabcolsep}{3pt}
\begin{tabular}{lccccc}
\toprule
Environment & $1-q_0$ & $q_0$ & $k_{\rm cross}$ [m$^{-1}$] & $\lambda_{\rm cross}$ [m] & $f_{\rm cross}$ [Hz] \\
\midrule
IGM & $10^{-7}$ & $\approx 1$ & $3.1{\times}10^{27}$ & $2.0{\times}10^{-27}$ & $1.5{\times}10^{35}$ \\
Solar surface & $2{\times}10^{-6}$ & $\approx 1$ & $6.2{\times}10^{28}$ & $1.0{\times}10^{-28}$ & $3.0{\times}10^{36}$ \\
Earth surface & $7{\times}10^{-10}$ & $\approx 1$ & $2.2{\times}10^{25}$ & $2.9{\times}10^{-25}$ & $1.0{\times}10^{33}$ \\
Laboratory & $10^{-27}$ & $\approx 1$ & $3.1{\times}10^{7}$ & $2.0{\times}10^{-7}$ & $1.5{\times}10^{15}$ \\
Near BH ($r{=}r_s$) & $0.393$ & $0.607$ & $9.5{\times}10^{33}$ & $6.6{\times}10^{-34}$ & $4.5{\times}10^{41}$ \\
Max.\ crossover ($q_0{=}1/3$) & $2/3$ & $1/3$ & $1.2{\times}10^{34}$ & $5.3{\times}10^{-34}$ & $5.7{\times}10^{41}$ \\
\bottomrule
\end{tabular}
\end{table}

\textbf{Key observations:}
\begin{itemize}
\item In all astrophysical and laboratory environments with $q_0\approx 1$, the crossover wavelength is sub-Planckian ($\lambda_{\rm cross} \ll \ell_P$). All physically accessible modes are in the wave regime.
\item Only near saturation ($q_0 \ll 1$) does the crossover enter potentially observable scales. However, $q_0 \ll 1$ is the regime where the continuum description breaks down.
\item The practical upshot: \textbf{all observable gravitational waves propagate in the wave regime} at speed $c$. No frequency-dependent dispersion is observable with current or planned detectors.
\end{itemize}

% ============================================================
\subsection{D4. GW170817 Analysis --- Integrated Phase Velocity Through the Milky Way Halo}
% ============================================================

\subsubsection{D4.1 The GW170817 constraint}

GW170817 constrained $|v_{\rm GW} - c|/c \le 10^{-15}$. We compute the integrated phase velocity correction through the Milky Way halo to determine whether the framework satisfies this bound.

\subsubsection{D4.2 Phase velocity in a $q$-field background}

For a wave propagating through a region with local $q_0(\mathbf{x})$, the phase velocity is (from the dispersion relation in the wave regime):
\begin{equation}
v_\phi(\mathbf{x}) = \frac{c}{\sqrt{q_0(\mathbf{x})}} \approx c\left(1 + \frac{1}{2}(1-q_0(\mathbf{x}))\right),
\end{equation}
where we expanded for $1-q_0 \ll 1$ (valid everywhere except very near compact objects).

\subsubsection{D4.3 Milky Way halo model}

The Milky Way halo has mass $M_{\rm halo} \sim 10^{12} M_\odot$, extending to $R_{\rm halo} \sim 200$~kpc. The $q$-field at a distance $r$ from the galactic center is:
\begin{equation}
q_0(r) \approx e^{-GM_{\rm halo}/(r c^2)} \approx 1 - \frac{GM_{\rm halo}}{r c^2}.
\end{equation}

At the solar circle ($r \approx 8$~kpc):
\begin{equation}
1 - q_0(R_\odot^{\rm gal}) \approx \frac{GM_{\rm halo}}{R_\odot^{\rm gal} c^2} \sim \frac{6.67\times 10^{-11}\times 10^{12}\times 2\times 10^{30}}{8\times 3.09\times 10^{19}\times 9\times 10^{16}} \approx 6\times 10^{-6}.
\end{equation}

\subsubsection{D4.4 Integrated phase delay}

The GW170817 signal traveled $\sim 40$~Mpc from NGC 4993 to Earth. The dominant $q$-field gradient is near the source and in the Milky Way. Integrating along the line of sight:
\begin{equation}
\frac{\Delta v_\phi}{c} \sim \frac{1}{2}\int_{\rm LOS} (1-q_0(\mathbf{x}))\,\frac{d\ell}{D} \sim \frac{GM_{\rm MW}}{2D c^2} \sim 6\times 10^{-10},
\label{eq:gw170817_correction}
\end{equation}
where we used the fact that the integral of $(1-q_0)$ along the line of sight scales as $GM/D c^2$, and the total mass along the path (integrated) is dominated by the Milky Way. The factor of $1/D$ (where $D \sim 40$~Mpc is the total distance) suppresses the effect by $\sim 10^{-7}$ relative to the local $(1-q_0) \sim 10^{-6}$ at the solar circle.

\textbf{A more careful computation:} The cumulative phase shift is:
\begin{equation}
\Delta\phi = \omega\int_0^D \left(\frac{1}{v_\phi(\ell)} - \frac{1}{c}\right)d\ell \approx -\frac{\omega}{2c}\int_0^D (1-q_0(\ell))\,d\ell.
\end{equation}

The integral $\int (1-q_0(\ell))d\ell$ is dominated by the near-field contribution from the host galaxy (within $\sim 100$~kpc of the source) and the Milky Way (within $\sim 200$~kpc of the detector). The intergalactic contribution is negligible ($1-q_0 \sim 10^{-7}$ over $\sim 40$~Mpc gives $\sim 4\times 10^{-7}$~Mpc $\sim 1.2\times 10^{16}$~m, times $(1-q_0) \sim 10^{-7}$, gives $\sim 10^9$~m integrated deficit).

Using the exponential solution $q_0(r)=e^{-GM/rc^2}$ for the Milky Way:
\begin{equation}
\int_{\rm MW} (1-q_0(\ell))\,d\ell \approx \frac{GM_{\rm MW}}{c^2}\ln\frac{R_{\rm halo}}{R_{\rm min}} \sim \frac{GM_{\rm MW}}{c^2}\times \mathcal{O}(10) \sim 1.5\times 10^{16}~\text{m}.
\end{equation}

With $\omega \sim 2\pi\times 100$~Hz $\sim 600$~s$^{-1}$:
\begin{equation}
\Delta\phi \sim \frac{600}{2\times 3\times 10^8}\times 1.5\times 10^{16} \sim 1.5\times 10^{10}~\text{radians}.
\end{equation}

This seems large, but it's a \emph{static} phase shift (the $q$-field is static on GW timescales). The observable is the \emph{frequency-dependent} part of the phase. Since $1-q_0$ varies on galactic scales ($\sim$~kpc), while GW wavelengths are $\sim 10^3$~km, the static phase shift is constant across the LIGO band and does not produce a measurable time delay.

The \textbf{differential} speed between two frequencies $f_1, f_2$ is:
\begin{equation}
\frac{v_\phi(f_1) - v_\phi(f_2)}{c} \approx \frac{1}{2}(1-q_0)\left[g(k_1) - g(k_2)\right],
\end{equation}
where $g(k)$ is the dispersive correction from finite $\gamma_0(1-q_0)$. At LIGO frequencies, $k \ll k_{\rm cross}$ by $\sim 40$ orders of magnitude, so $g(k_1)-g(k_2)$ is utterly negligible. The dispersive correction is:
\begin{equation}
\frac{\Delta v_\phi}{c}\bigg|_{\rm disp} \sim \left(\frac{k}{k_{\rm cross}}\right)^2 \sim 10^{-80}\ \text{--}\ 10^{-70}.
\end{equation}

\subsubsection{D4.5 Conclusion}

\textbf{Consistency with GW170817.} The framework predicts $v_{\rm GW} = c$ in vacuum to within one part in $10^{70}$ (non-dispersive static shift) or $10^{80}$ (dispersive correction)---both astronomically below the $10^{-15}$ constraint. However, as noted in the main paper (Sec.~8.3), this analysis covers scalar $q$-field modes; the observed GW170817 signal may correspond to tensor metric perturbations whose decoupling from the scalar sector is conjectured but not yet derived. Full consistency therefore requires proof of tensor-scalar decoupling at linear order, which remains an open problem.

% ============================================================
\subsection{D5. Frequency-Dependent GW Propagation --- Predicted Deviations from GR at Low Frequencies}
% ============================================================

\subsubsection{D5.1 The dispersive regime}

From the full dispersion relation (\ref{eq:full_dispersion}), frequency-dependent effects become significant when $k \sim k_{\rm cross}$, i.e., when:
\begin{equation}
f \sim f_{\rm cross} \equiv \frac{c k_{\rm cross}}{2\pi} = \frac{\gamma_0(1-q_0)}{4\pi}\sqrt{q_0}.
\end{equation}

In vacuum ($q_0 \to 1$): $f_{\rm cross} \to 0$. \textbf{There is no low-frequency cutoff in vacuum.} All frequencies propagate as waves.

At finite $(1-q_0)$: $f_{\rm cross} > 0$. Modes with $f < f_{\rm cross}$ are overdamped and decay without oscillation.

\subsubsection{D5.2 Where could the crossover be observable?}

For $f_{\rm cross}$ to enter the LIGO band ($f \sim 10$--$1000$~Hz), we need:
\begin{equation}
\gamma_0(1-q_0) \sim 4\pi f \sim 10^2\text{--}10^4~\text{s}^{-1} \quad\Longrightarrow\quad (1-q_0) \sim 10^{-41}\text{--}10^{-39}.
\end{equation}

This requires $1-q_0 \sim 10^{-40}$, which corresponds to the $q$-field at a distance:
\begin{equation}
r = \frac{GM}{c^2(1-q_0)} \sim \frac{10^3}{10^{-40}} = 10^{43}~\text{m} \sim 10^{27}~\text{Mpc},
\end{equation}
far beyond the observable universe. \textbf{The crossover is not observable with current or planned detectors.}

\subsubsection{D5.3 What could change this}

If the cell spacing $a$ is significantly larger than $\ell_P$, then $\gamma_0 = c/a$ is smaller, and $k_{\rm cross}$ shifts to lower (more accessible) values. For $a \sim 10^{-15}$~m (nuclear scale, as in the now-corrected $N\sim 10^{90}$ estimate), $\gamma_0 \sim 3\times 10^{23}$~s$^{-1}$, and $f_{\rm cross}$ in the solar neighborhood would be $\sim 10^{-14}$~Hz --- within the pulsar timing array band. However, as shown in \S F3, the nuclear-scale identification is inconsistent with the Newtonian calibration of $\kappa$.

\subsubsection{D5.4 The honest conclusion}

With $\gamma_0$ fixed by the Planck scale and $c$ fixed by observation, the wave-diffusion crossover is a Planck-scale phenomenon. Its significance is theoretical, not observational: it provides a conceptual unification of wave (coherent) and diffusive (decoherent) gravitational dynamics within a single equation, with the crossover controlled by the capacity-dependent damping $\gamma_0(1-q)$.

\textbf{Flag:} All predictions in D1-D5 depend on the unified telegraph equation (\ref{eq:a3_telegraph}) and the coupling to matter via (\ref{eq:scalar_wave}). There are no ``scalar GW polarization'' predictions. The framework's sole clean GW prediction is the vacuum wave speed $c$, which is exact and consistent with GW170817.

%%%%%%%%%%%%%%%%%%%%%%%%%%%%%%%%%%%%%%%%%%%%%%%%%%%%%%%%%%%%%%%%%%%%%
\section{Inspector Audit Trail (v3)}
%%%%%%%%%%%%%%%%%%%%%%%%%%%%%%%%%%%%%%%%%%%%%%%%%%%%%%%%%%%%%%%%%%%%%

This section catalogues every issue found by internal adversarial review, with resolution status.

% ============================================================
\subsection{E1. Complete Issue Inventory with V2/V3 Resolution Status}
% ============================================================

\begin{table}[h]
\caption{Complete audit issue inventory with resolution status.\\
RESOLVED = mathematical proof provided; RESOLVED-BY-REMOVAL = claim/feature removed; RESOLVED-BY-REWRITE = section rewritten; OPEN = under investigation.}
\begin{tabular}{p{0.5cm}p{5.5cm}ll}
\toprule
ID & Issue & Severity & V2/V3 Status \\
\midrule
G1 & Independence circularity in Shannon derivation & FATAL & RESOLVED-BY-REWRITE \\
G6 & Uniqueness-vs-universality trilemma & FATAL & RESOLVED-BY-REWRITE \\
H2 & Shadow prediction from ad hoc refractive index & FATAL & RESOLVED-BY-REMOVAL \\
H6 & EHT M87* excludes DGF shadow at 7$\sigma$ & FATAL & RESOLVED-BY-REMOVAL \\
H1 & Parabolic equation cannot describe GWs & FATAL & RESOLVED-BY-REWRITE \\
F1 & \textbf{NEW:} $\gamma$ was a free parameter, not derived & FATAL (new) & RESOLVED \\
F2 & \textbf{NEW:} No preferred-frame containment argument & SEVERE (new) & RESOLVED \\
F3 & \textbf{NEW:} Lorentz emergence not rigorous & SEVERE (new) & RESOLVED \\
F4 & \textbf{NEW:} Wave speed $c_s$ required separate calibration & SEVERE (new) & RESOLVED-BY-REWRITE \\
F5 & \textbf{NEW:} Telegraph equation had constant $\gamma$, inconsistent with capacity & SEVERE (new) & RESOLVED \\
F6 & \textbf{NEW:} Lyapunov functional only for constant $\gamma$ & MODERATE (new) & RESOLVED \\
F7 & \textbf{NEW:} $N\sim 10^{90}$ inconsistent with Planck calibration & MODERATE (new) & RESOLVED-BY-REWRITE \\
G5b & Shannon citation inflation & SEVERE & RESOLVED-BY-REWRITE \\
G2 & Bridge from Shannon to flux: postulate, not derivation & SEVERE & RESOLVED-BY-REWRITE \\
G4 & Irrotationality and gradient flow not in assumptions & SEVERE & RESOLVED-BY-REWRITE \\
G5a & B1 not acknowledged as assumption & SEVERE & RESOLVED-BY-REWRITE \\
H4 & Ringdown estimate excluded by LIGO & SEVERE & RESOLVED-BY-REMOVAL \\
H5 & Information paradox ``resolution'' relabeling & SEVERE & RESOLVED-BY-REWRITE \\
A8 & Entropy can decrease between steps & SEVERE & RESOLVED-BY-REWRITE \\
I1 & Conflation of ``consistent with'' and ``forced by'' & SEVERE & RESOLVED-BY-REWRITE \\
H3 & Physical state at $q\to 0$ cores undefined & SEVERE & OPEN \\
M1 & \textbf{NEW:} Audit trail contained self-congratulatory language & MODERATE (new) & RESOLVED-BY-REWRITE \\
M2 & \textbf{NEW:} Scalar GW polarization (D2-D4) not retracted clearly & SEVERE (new) & RESOLVED-BY-REMOVAL \\
M3 & \textbf{NEW:} Damping derivation was post-hoc, not from Definition~C & SEVERE (new) & RESOLVED \\
M4 & \textbf{NEW:} $c_s$ calibration was an independent postulate & MODERATE (new) & RESOLVED-BY-REWRITE \\
M5 & \textbf{NEW:} No proof that boundary terms vanish in Lyapunov & MINOR (new) & RESOLVED \\
M6 & \textbf{NEW:} $\gamma\to\infty$ singular perturbation obscures physics & MINOR (new) & RESOLVED-BY-REWRITE \\
C4 & Sandpile analogy not quantitatively accurate & MODERATE & OPEN \\
G3 & Choice of vacant vs.\ occupied self-information & MODERATE & RESOLVED-BY-REWRITE \\
G7 & C5 (vacuum stability) is a postulate, not a definition & MODERATE & RESOLVED-BY-REWRITE \\
B8 & No quantitative prediction that is (a) derived, (b) not excluded & MODERATE & OPEN \\
D8 & Scale-dependence title contradiction & MODERATE & RESOLVED-BY-REWRITE \\
E1 & Flux function chosen, not derived & MODERATE & RESOLVED-BY-REWRITE \\
C6 & Inflation explanation requires $q_i$ plausibility & MINOR & OPEN \\
\bottomrule
\end{tabular}
\end{table}

Of the 13 new issues found by earlier adversarial review, all are resolved. F1 ($\gamma$ was a free parameter) was the most severe: it was resolved by motivating $\gamma(q)=\gamma_0(1-q)$ from Definition~C (the linear form is flagged as a modeling ansatz in the main paper). F2 and F3 fell when we showed the damping term is second-order in $\delta q$ around $q=1$, so Lorentz invariance---always present as the d'Alembertian's symmetry---is restored in vacuum. Four problems remain open and deserve attention in any future extension.

% ============================================================
\subsection{E2. Resolution Details for New Issues}
% ============================================================

\begin{enumerate}[leftmargin=*]
\item[\textbf{F1:}] \textbf{$\gamma$ was a free parameter.} \textit{Resolution:} Motivated $\gamma(q)=\gamma_0(1-q)$ from Definition~C (occupied cells block causal events with probability $1-q$). The linear form is the simplest ansatz; the essential property $\gamma\to 0$ as $q\to 1$ follows from the capacity constraint. Flagged as a modeling choice in the main paper (Sec.~IV). See \S A3.1.

\item[\textbf{F2:}] \textbf{No preferred-frame containment.} \textit{Resolution:} Showed that $\gamma(q)=0$ exactly at $q=1$ (vacuum). Preferred-frame effects, which require $\gamma>0$, are therefore confined to regions with $q<1$, i.e., near mass. The suppression factor is $(1-q) \sim GM/rc^2 \sim 10^{-6}$ at the solar surface (corrected from an earlier estimate). Lorentz symmetry with respect to $c$ is a structural feature of the d'Alembertian, not emergent; the vacuum restores this pre-existing symmetry. See \S A3.4.

\item[\textbf{F3:}] \textbf{Lorentz emergence not rigorous.} \textit{Resolution:} Showed that damping $\gamma_0(1-q)\partial_t q$ is $\mathcal{O}(\delta q^2)$ around $q=1$, so the linearized dynamics are exactly $\Box(\delta q)=0$. Lorentz invariance of the wave operator was always present (the d'Alembertian with speed $c$ is built into the current equation); vacuum restores this symmetry by suppressing the damping term. See \S A3.4 and main paper Sec.~V.

\item[\textbf{F4:}] \textbf{Wave speed required calibration.} \textit{Resolution:} In the unified framework, the wave speed is $c$ by construction (from the $c^2/a^2$ coupling in the current equation). No separate calibration $\kappa/\tau = c^2 q_0$ is needed. See \S A3.7.

\item[\textbf{F5:}] \textbf{Telegraph equation had constant $\gamma$.} \textit{Resolution:} The telegraph equation has capacity-dependent damping $\gamma_0(1-q)$. The constant-$\gamma$ case is recovered only as the leading approximation when $q$ is nearly uniform; the full nonlinear equation includes the $(1-q)$ factor. See \S A3.3.

\item \textbf{F6: Lyapunov functional only for constant $\gamma$.} \textit{Resolution:} Proved $dE/dt \le 0$ for the generalized energy functional with capacity-dependent damping $\gamma_0(1-q)$. See \S A7.

\item \textbf{F7: $N\sim 10^{90}$ inconsistent with Planck calibration.} \textit{Resolution:} Corrected to $N\sim 10^{184}$ (Planck-scale cells). The $N\sim 10^{90}$ estimate used nuclear-scale cells ($\sim 10^{-15}$~m), which are inconsistent with the Newtonian calibration requiring $\kappa \sim \ell_P c$. See \S B5, \S F3.
\end{enumerate}

% ============================================================
\subsection{E3. Remaining Open Problems}
% ============================================================

\begin{enumerate}[leftmargin=*, label=\textbf{OP-\arabic*:}]

\item \textbf{Cell basis problem.} What selects the $\{|0\rangle,|1\rangle\}$ basis? \textbf{Status: Unproved conjecture.} Until resolved, the $q$-field is defined relative to an externally specified preferred basis.

\item \textbf{Nonlinear damping $\gamma(q)$.} The linear ansatz $\gamma(q)=\gamma_0(1-q)$ is the simplest form. Nonlinear corrections may exist. \textbf{Status: Modeling ansatz.} Core consequence ($\gamma\to 0$ as $q\to 1$) is independent of specific form.

\item \textbf{Coupling to matter and radiation.} How does the $q$-field couple to matter? Is there an effective metric $g_{\mu\nu}[q]$? \textbf{Status: Completely open.}

\item \textbf{Strong-field regime.} The exact solution to $\nabla^2(\ln q)=0$ with boundary condition $q(r_s)\to 0$ is not known in closed form for extended sources. \textbf{Status: Conjectured but not proved.}

\item \textbf{Spatial dimensionality.} DGF does not predict $d=3$. \textbf{Status: $d=3$ is an observational input.}

\item \textbf{Numerical constants: $G$, $c$, $\hbar$.} The coupling $4\pi G/c^2$ is calibrated, not derived. \textbf{Status: Acknowledged calibration.}

\item \textbf{Born rule.} Not derived from Definition~C. \textbf{Status: Not addressed.} (Zilly's theorem provides the Born rule from finite capacity plus cyclic dynamics; we accept it as input.)

\item \textbf{Correlation corrections to the flux.} The self-information flux assumes uncorrelated cells. \textbf{Status: Open.}

\item \textbf{Discrete-to-continuum limit for the full dynamics.} Deriving the unified PDE from the discrete permutation dynamics is an open combinatorial problem. \textbf{Status: Not addressed.}

\item \textbf{Cosmology.} DGF has no cosmological extension. \textbf{Status: Open research direction.}

\end{enumerate}

\textbf{Honest assessment:} Of the 10 open problems, OP-3 (coupling to matter) is the most urgent for empirical contact beyond the static Newtonian correspondence. OP-1 (cell basis) is the deepest conceptual issue. These two are the ones a referee is most likely to ask about.

%%%%%%%%%%%%%%%%%%%%%%%%%%%%%%%%%%%%%%%%%%%%%%%%%%%%%%%%%%%%%%%%%%%%%
\section{Numerical Estimates}
%%%%%%%%%%%%%%%%%%%%%%%%%%%%%%%%%%%%%%%%%%%%%%%%%%%%%%%%%%%%%%%%%%%%%

This section collects all numerical values used in the paper, with their derivations. No number appears without justification.

% ============================================================
\subsection{F1. $\gamma_0$ from Planck-Scale Identification}
% ============================================================

\subsubsection{F1.1 Definition}

The damping scale $\gamma_0 = c/a$. With the cell spacing identified as the Planck length $a = \ell_P = \sqrt{\hbar G/c^3} \approx 1.616\times 10^{-35}$~m:
\begin{equation}
\boxed{\gamma_0 = \frac{c}{\ell_P} = \frac{1}{t_P} \approx 1.855\times 10^{43}~\text{s}^{-1}}.
\label{eq:gamma0_value}
\end{equation}

This is the Planck frequency. The corresponding timescale is the Planck time $t_P \approx 5.391\times 10^{-44}$~s.

\subsubsection{F1.2 Consistency with Newtonian calibration}

The static calibration requires the source term $\nabla^2\psi = -(4\pi G/c^2)\rho_m$ to reproduce Newtonian gravity. In the DGF framework, $G$ is an empirical input (Tier~1); the coupling constant linking the $q$-field to matter is not derived from $\gamma_0$ or $\ell_P$, but calibrated to the observed value of Newton's constant. The $\mathcal{O}(1)$ dimensionless prefactor in the coupling is a free parameter of the theory, to be constrained by observation.

\textbf{Bottom line:} $\gamma_0$ is fixed by known constants ($c$, $\ell_P$). The theory has \textbf{zero additional free parameters} for the damping scale.

% ============================================================
\subsection{F2. Diffusion-to-Wave Crossover --- Numerical Values}
% ============================================================

From \S D2, the crossover wavenumber is:
\begin{equation}
k_{\rm cross} = \frac{\gamma_0(1-q_0)}{2c}\sqrt{q_0}.
\end{equation}

For $q_0 \approx 1$ (all realistic environments):
\begin{equation}
k_{\rm cross} \approx \frac{\gamma_0(1-q_0)}{2c}.
\end{equation}

\textbf{Intergalactic medium} ($1-q \sim 10^{-7}$): $k_{\rm cross} \sim 3\times 10^{27}$~m$^{-1}$, $\lambda_{\rm cross} \sim 2\times 10^{-27}$~m $\sim 10^{8}\,\ell_P$.

\textbf{Solar surface} ($1-q \sim 2\times 10^{-6}$): $k_{\rm cross} \sim 6\times 10^{28}$~m$^{-1}$, $\lambda_{\rm cross} \sim 1\times 10^{-28}$~m $\sim 6\times 10^{6}\,\ell_P$.

\textbf{Laboratory} ($1-q \sim 10^{-27}$): $k_{\rm cross} \sim 3\times 10^{7}$~m$^{-1}$, $\lambda_{\rm cross} \sim 2\times 10^{-7}$~m (optical wavelength).

\textbf{Implication:} In all environments except extreme saturation ($q_0 \ll 1$), the crossover wavelength is below or near the Planck scale. All observable gravitational phenomena are in the wave regime.

% ============================================================
\subsection{F3. Planck-Scale Parameters with Updated Notation}
% ============================================================

\begin{center}
\begin{tabular}{lll}
\toprule
Parameter & Symbol & Value/Relation \\
\midrule
Cell spacing & $a$ & $\ell_P = 1.616\times 10^{-35}$~m (input) \\
Fundamental speed & $c$ & $2.998\times 10^8$~m/s (input) \\
Damping scale & $\gamma_0$ & $c/a = 1/t_P \approx 1.855\times 10^{43}$~s$^{-1}$ \\
Coupling constant & $c^2/a^2$ & $\approx 3.45\times 10^{86}$~m$^2$/s$^2$ (derived) \\
Cell count & $N$ & $(c/H_0\ell_P)^3 \sim 10^{184}$ (fiducial) \\
Gravity coupling & $4\pi G/c^2$ & $9.33\times 10^{-27}$~m$^2$/s (calibrated) \\
\bottomrule
\end{tabular}
\end{center}

\textbf{Key relationships:}
\begin{itemize}
\item $\gamma_0 = c/a$: the natural frequency scale is the inverse light-crossing time between adjacent cells.
\item $c^2/a^2 = c^2/\ell_P^2$: the coupling in the unified current equation is set by the Planck area.
\item $G$ is not predicted; it is calibrated via the static Poisson correspondence. The relation $G = \alpha\ell_P c^3/(4\pi)$ with $\alpha \approx 2$ links $G$ to $\ell_P$ and $c$, but the $\mathcal{O}(1)$ factor $\alpha$ is not derived.
\end{itemize}

%%%%%%%%%%%%%%%%%%%%%%%%%%%%%%%%%%%%%%%%%%%%%%%%%%%%%%%%%%%%%%%%%%%%%
\section{Notation Table: SM $\leftrightarrow$ Paper Parameter Mapping}
%%%%%%%%%%%%%%%%%%%%%%%%%%%%%%%%%%%%%%%%%%%%%%%%%%%%%%%%%%%%%%%%%%%%%

A complete mapping between SM notation, main paper notation, and deprecated notation. Nothing should be ambiguous when cross-referencing.

\begin{table}[h]
\caption{Notation mapping: paper $\leftrightarrow$ this SM $\leftrightarrow$ deprecated.}
\label{tab:notation}
\begin{tabular}{llll}
\toprule
Quantity & Paper \& This SM & Deprecated & Notes \\
\midrule
Vacant fraction & $q$ & $q$ & unchanged \\
Current (edge) & $J_{ij}$ & $J_{i\to j}$ & simplified notation \\
Current (continuum) & $\mathbf{J}$ & $\mathbf{J}$ & unchanged \\
Fundamental speed & $c$ & $c$, $c_s$ & $c_s$ was calibrated; now $c$ is fundamental \\
Cell spacing & $a$ & $a$ & unchanged; $a=\ell_P$ \\
Damping scale & $\gamma_0 = c/a$ & $\gamma = 1/\tau$ & $\gamma_0$ is capacity-independent base rate \\
Effective damping & $\gamma_0(1-q)$ & $\gamma$ (constant) & key innovation \\
Coupling constant & $c^2/a^2$ & $\kappa$, $\kappa/\tau$ & $\kappa$ was diffusivity; coupling now split \\
Edge-averaged $q$ & $\bar{q} = (q_i+q_j)/2$ & $q_0$ (background) & $\bar{q}$ is edge-local; $q_0$ was global \\
Information potential & $\psi = -\ln q$ & $\psi = -\ln q$ & unchanged \\
Free energy density & $\Phi(q)=q\ln q - q + 1$ & $\Phi(q)=q\ln q - q + 1$ & unchanged \\
D'Alembertian & $\Box = \partial_t^2 - c^2\nabla^2$ & not used previously & new \\
Crossover wavenumber & $k_{\rm cross}$ & $k_{\rm cross}$ & formula updated with $q$-dependence \\
Singular perturbation & $\varepsilon = \omega/[\gamma_0(1-q)]$ & ``$\gamma\to\infty$'' & physically meaningful parameter \\
Cell count & $N \sim 10^{184}$ & $N \sim 10^{90}$ & corrected (see \S B5) \\
\bottomrule
\end{tabular}
\end{table}

\textbf{Critical rule:} The coupling constant is $\kappa\equiv c^2/a^2$ (used interchangeably in the main paper's H-theorem section). The damping is $\gamma_0(1-q)$, NOT $1/\tau$ or constant $\gamma$. Any SM equation using $\tau$ or constant $\gamma$ is referencing deprecated notation and is inconsistent with the present paper.

%%%%%%%%%%%%%%%%%%%%%%%%%%%%%%%%%%%%%%%%%%%%%%%%%%%%%%%%%%%%%%%%%%%%%
\section*{References for Supplemental Material}
%%%%%%%%%%%%%%%%%%%%%%%%%%%%%%%%%%%%%%%%%%%%%%%%%%%%%%%%%%%%%%%%%%%%%

\begin{thebibliography}{99}

\bibitem{Cattaneo1948} C.~Cattaneo, Atti Sem.\ Mat.\ Fis.\ Univ.\ Modena \textbf{3}, 83 (1948).
\bibitem{Vernotte1958} P.~Vernotte, C.\ R.\ Acad.\ Sci.\ \textbf{246}, 3154 (1958).
\bibitem{Joseph1989} D.D.\ Joseph and L.\ Preziosi, Rev.\ Mod.\ Phys.\ \textbf{61}, 41 (1989).
\bibitem{Jackson1970} H.E.\ Jackson and C.T.\ Walker, ``Second Sound in NaF,'' Phys.\ Rev.\ Lett.\ \textbf{25}, 26--28 (1970).
\bibitem{London1935} F.\ London and H.\ London, Proc.\ R.\ Soc.\ A \textbf{149}, 71 (1935).
\bibitem{Maxwell1865} J.C.\ Maxwell, Phil.\ Trans.\ R.\ Soc.\ \textbf{155}, 459 (1865).
\bibitem{Heaviside1876} O.\ Heaviside, Phil.\ Mag.\ \textbf{2}, 135 (1876).
\bibitem{Goldstein1951} S.~Goldstein, Quart.\ J.\ Mech.\ Appl.\ Math.\ \textbf{4}, 129 (1951).
\bibitem{Weiss1994} G.H.\ Weiss, \emph{Aspects and Applications of the Random Walk} (North-Holland, 1994).
\bibitem{Caldeira1983} A.O.\ Caldeira and A.J.\ Leggett, Ann.\ Phys.\ \textbf{149}, 374 (1983).
\bibitem{Will2014} C.M.\ Will, Living Rev.\ Relativ.\ \textbf{17}, 4 (2014).
\bibitem{LIGO2017} B.P.\ Abbott et al.\ (LIGO/Virgo), Astrophys.\ J.\ \textbf{848}, L13 (2017).
\bibitem{LIGO2019} B.P.\ Abbott et al.\ (LIGO/Virgo), Phys.\ Rev.\ Lett.\ \textbf{123}, 011102 (2019).
\bibitem{EHT2019} K.~Akiyama et al.\ (EHT), Astrophys.\ J.\ \textbf{875}, L1 (2019).
\bibitem{Shannon1948} C.E.\ Shannon, Bell Syst.\ Tech.\ J.\ \textbf{27}, 379 (1948).
\bibitem{Vazquez2007} J.L.\ Vazquez, \emph{The Porous Medium Equation} (Oxford, 2007).
\bibitem{Bak1987} P.~Bak, C.~Tang, and K.~Wiesenfeld, Phys.\ Rev.\ Lett.\ \textbf{59}, 381 (1987).
\bibitem{Zurek2009} W.H.\ Zurek, Nat.\ Phys.\ \textbf{5}, 181 (2009).
\end{thebibliography}

\end{document}
